%% Supplementary Information for
%% "Which metric should choose the model? Selection criteria, class weighting and
%% endpoint recall in tabular ordinal classification".
%% Standard article class; compiles with any LaTeX distribution. Tables are generated
%% by analysis/make_tables.py.
\documentclass[10pt]{article}
\usepackage[a4paper,margin=2cm]{geometry}
\usepackage{booktabs,longtable,array}
\usepackage{amsmath,amssymb}
\usepackage[hidelinks]{hyperref}
\renewcommand{\thetable}{S\arabic{table}}
\setlength{\LTcapwidth}{\textwidth}
% Generated by analysis/make_tables.py from analysis/out; do not edit by hand.
\newcommand{\NDecisions}{72}
\newcommand{\NBoot}{2,000}
\newcommand{\NRecompute}{1,120}
\newcommand{\RecomputeMaxDiff}{\ensuremath{1.7\times 10^{-16}}}
\newcommand{\BalQwk}{0.72}
\newcommand{\RegEndQwk}{0.55}
\newcommand{\RegQwkQwk}{0.06}
\newcommand{\RegMaeQwk}{0.48}
\newcommand{\RegMmaeQwk}{0.29}
\newcommand{\RankEndQwk}{4.08}
\newcommand{\RankQwkQwk}{2.25}
\newcommand{\AgrQwkQwk}{1.00}
\newcommand{\AgrQwkMae}{0.22}
\newcommand{\AgrQwkAmae}{0.53}
\newcommand{\AgrQwkMmae}{0.47}
\newcommand{\AgrQwkMfone}{0.79}
\newcommand{\AgrQwkRps}{0.18}
\newcommand{\AgrQwkCon}{0.56}
\newcommand{\BalMae}{0.00}
\newcommand{\RegEndMae}{0.89}
\newcommand{\RegQwkMae}{0.55}
\newcommand{\RegMaeMae}{0.01}
\newcommand{\RegMmaeMae}{0.91}
\newcommand{\RankEndMae}{5.58}
\newcommand{\RankQwkMae}{4.75}
\newcommand{\AgrMaeQwk}{0.22}
\newcommand{\AgrMaeMae}{1.00}
\newcommand{\AgrMaeAmae}{0.19}
\newcommand{\AgrMaeMmae}{0.01}
\newcommand{\AgrMaeMfone}{0.18}
\newcommand{\AgrMaeRps}{0.68}
\newcommand{\AgrMaeCon}{0.00}
\newcommand{\BalAmae}{0.68}
\newcommand{\RegEndAmae}{0.48}
\newcommand{\RegQwkAmae}{0.35}
\newcommand{\RegMaeAmae}{0.56}
\newcommand{\RegMmaeAmae}{0.42}
\newcommand{\RankEndAmae}{4.04}
\newcommand{\RankQwkAmae}{3.92}
\newcommand{\AgrAmaeQwk}{0.53}
\newcommand{\AgrAmaeMae}{0.19}
\newcommand{\AgrAmaeAmae}{1.00}
\newcommand{\AgrAmaeMmae}{0.43}
\newcommand{\AgrAmaeMfone}{0.51}
\newcommand{\AgrAmaeRps}{0.14}
\newcommand{\AgrAmaeCon}{0.57}
\newcommand{\BalMmae}{0.93}
\newcommand{\RegEndMmae}{0.29}
\newcommand{\RegQwkMmae}{0.36}
\newcommand{\RegMaeMmae}{0.78}
\newcommand{\RegMmaeMmae}{0.06}
\newcommand{\RankEndMmae}{2.29}
\newcommand{\RankQwkMmae}{4.79}
\newcommand{\AgrMmaeQwk}{0.47}
\newcommand{\AgrMmaeMae}{0.01}
\newcommand{\AgrMmaeAmae}{0.43}
\newcommand{\AgrMmaeMmae}{1.00}
\newcommand{\AgrMmaeMfone}{0.32}
\newcommand{\AgrMmaeRps}{0.01}
\newcommand{\AgrMmaeCon}{0.67}
\newcommand{\BalMfone}{0.72}
\newcommand{\RegEndMfone}{0.52}
\newcommand{\RegQwkMfone}{0.07}
\newcommand{\RegMaeMfone}{0.50}
\newcommand{\RegMmaeMfone}{0.27}
\newcommand{\RankEndMfone}{4.04}
\newcommand{\RankQwkMfone}{3.04}
\newcommand{\AgrMfoneQwk}{0.79}
\newcommand{\AgrMfoneMae}{0.18}
\newcommand{\AgrMfoneAmae}{0.51}
\newcommand{\AgrMfoneMmae}{0.32}
\newcommand{\AgrMfoneMfone}{1.00}
\newcommand{\AgrMfoneRps}{0.15}
\newcommand{\AgrMfoneCon}{0.43}
\newcommand{\BalRps}{0.00}
\newcommand{\RegEndRps}{0.81}
\newcommand{\RegQwkRps}{0.56}
\newcommand{\RegMaeRps}{0.02}
\newcommand{\RegMmaeRps}{0.91}
\newcommand{\RankEndRps}{5.38}
\newcommand{\RankQwkRps}{4.75}
\newcommand{\AgrRpsQwk}{0.18}
\newcommand{\AgrRpsMae}{0.68}
\newcommand{\AgrRpsAmae}{0.14}
\newcommand{\AgrRpsMmae}{0.01}
\newcommand{\AgrRpsMfone}{0.15}
\newcommand{\AgrRpsRps}{1.00}
\newcommand{\AgrRpsCon}{0.00}
\newcommand{\BalCon}{0.99}
\newcommand{\RegEndCon}{0.25}
\newcommand{\RegQwkCon}{0.40}
\newcommand{\RegMaeCon}{0.85}
\newcommand{\RegMmaeCon}{0.09}
\newcommand{\RankEndCon}{2.58}
\newcommand{\RankQwkCon}{4.50}
\newcommand{\AgrConQwk}{0.56}
\newcommand{\AgrConMae}{0.00}
\newcommand{\AgrConAmae}{0.57}
\newcommand{\AgrConMmae}{0.67}
\newcommand{\AgrConMfone}{0.43}
\newcommand{\AgrConRps}{0.00}
\newcommand{\AgrConCon}{1.00}
\newcommand{\RegOwnMax}{0.11}
\newcommand{\FriedEndChi}{26.4}
\newcommand{\FriedEndP}{\ensuremath{<}0.001}
\newcommand{\FriedQwkChi}{16.3}
\newcommand{\FriedQwkP}{\ensuremath{=}0.012}
\newcommand{\NemenyiCD}{2.60}
\newcommand{\WilPAmae}{0.820}
\newcommand{\WilPosAmae}{3}
\newcommand{\WilNegAmae}{6}
\newcommand{\WilMedAmae}{\ensuremath{-}0.003}
\newcommand{\WilMeanAmae}{\ensuremath{+}0.004}
\newcommand{\QwkCostMedAmae}{\ensuremath{-}0.009}
\newcommand{\QwkCostNegAmae}{7}
\newcommand{\QwkCostPAmae}{0.039}
\newcommand{\WilPMfone}{0.938}
\newcommand{\WilPosMfone}{4}
\newcommand{\WilNegMfone}{3}
\newcommand{\WilMedMfone}{0.000}
\newcommand{\WilMeanMfone}{\ensuremath{+}0.004}
\newcommand{\QwkCostMedMfone}{0.000}
\newcommand{\QwkCostNegMfone}{7}
\newcommand{\QwkCostPMfone}{0.195}
\newcommand{\WilPCon}{0.031}
\newcommand{\WilPosCon}{6}
\newcommand{\WilNegCon}{0}
\newcommand{\WilMedCon}{\ensuremath{+}0.024}
\newcommand{\WilMeanCon}{\ensuremath{+}0.086}
\newcommand{\QwkCostMedCon}{\ensuremath{-}0.020}
\newcommand{\QwkCostNegCon}{6}
\newcommand{\QwkCostPCon}{0.031}
\newcommand{\WilPMmae}{0.039}
\newcommand{\WilPosMmae}{7}
\newcommand{\WilNegMmae}{1}
\newcommand{\WilMedMmae}{\ensuremath{+}0.002}
\newcommand{\WilMeanMmae}{\ensuremath{+}0.050}
\newcommand{\QwkCostMedMmae}{\ensuremath{-}0.010}
\newcommand{\QwkCostNegMmae}{8}
\newcommand{\QwkCostPMmae}{0.012}
\newcommand{\WilPMae}{0.055}
\newcommand{\WilPosMae}{2}
\newcommand{\WilNegMae}{7}
\newcommand{\WilMedMae}{\ensuremath{-}0.040}
\newcommand{\WilMeanMae}{\ensuremath{-}0.141}
\newcommand{\QwkCostMedMae}{\ensuremath{-}0.042}
\newcommand{\QwkCostNegMae}{8}
\newcommand{\QwkCostPMae}{0.020}
\newcommand{\WilPRps}{0.055}
\newcommand{\WilPosRps}{2}
\newcommand{\WilNegRps}{7}
\newcommand{\WilMedRps}{\ensuremath{-}0.034}
\newcommand{\WilMeanRps}{\ensuremath{-}0.129}
\newcommand{\QwkCostMedRps}{\ensuremath{-}0.039}
\newcommand{\QwkCostNegRps}{8}
\newcommand{\QwkCostPRps}{0.027}
\newcommand{\HolmCon}{0.094}
\newcommand{\AmaeMin}{\ensuremath{-}0.40}
\newcommand{\AmaeMax}{\ensuremath{+}0.53}
\newcommand{\ConQwkMedCell}{\ensuremath{-}0.020}
\newcommand{\ModRho}{\ensuremath{-}0.64}
\newcommand{\ModP}{0.026}
\newcommand{\ModCells}{12}
\newcommand{\CtrlPosEnd}{21}
\newcommand{\CtrlNegEnd}{3}
\newcommand{\CtrlMedEnd}{\ensuremath{+}0.145}
\newcommand{\CtrlPosMae}{24}
\newcommand{\CtrlNegMae}{0}
\newcommand{\CtrlMedMae}{\ensuremath{+}0.083}
\newcommand{\CtrlPosMmae}{2}
\newcommand{\CtrlNegMmae}{22}
\newcommand{\CtrlMedMmae}{\ensuremath{-}0.204}
\newcommand{\CtrlPosQwk}{17}
\newcommand{\CtrlNegQwk}{7}
\newcommand{\CtrlMedQwk}{\ensuremath{+}0.020}
\newcommand{\CtrlPosLow}{23}
\newcommand{\CtrlNegLow}{0}
\newcommand{\CtrlMedLow}{\ensuremath{+}0.320}
\newcommand{\CtrlPosHigh}{15}
\newcommand{\CtrlNegHigh}{8}
\newcommand{\CtrlMedHigh}{\ensuremath{+}0.121}
\newcommand{\CtrlPosAmae}{9}
\newcommand{\CtrlNegAmae}{15}
\newcommand{\CtrlMedAmae}{\ensuremath{-}0.020}
\newcommand{\CtrlN}{24}
\newcommand{\ConFallbacks}{14}
\newcommand{\ConFallbacksWine}{14}
\newcommand{\PredVarMlPoQwk}{2.68}
\newcommand{\EndShareMlPoQwk}{0.60}
\newcommand{\LowMlPoQwk}{0.52}
\newcommand{\HighMlPoQwk}{0.40}
\newcommand{\ObsMlPoQwk}{3.44}
\newcommand{\ExpMlPoQwk}{4.28}
\newcommand{\QwkMlPoQwk}{0.196}
\newcommand{\NamaeMlPoQwk}{0.354}
\newcommand{\NmmaeMlPoQwk}{0.388}
\newcommand{\PredVarMlPoAmae}{0.23}
\newcommand{\EndShareMlPoAmae}{0.06}
\newcommand{\LowMlPoAmae}{0.00}
\newcommand{\HighMlPoAmae}{0.11}
\newcommand{\ObsMlPoAmae}{1.41}
\newcommand{\ExpMlPoAmae}{1.63}
\newcommand{\QwkMlPoAmae}{0.137}
\newcommand{\NamaeMlPoAmae}{0.323}
\newcommand{\NmmaeMlPoAmae}{0.669}
\newcommand{\LabVarMl}{1.28}
\newcommand{\LabEndMl}{0.28}
\newcommand{\PredVarMlRfQwk}{1.98}
\newcommand{\EndShareMlRfQwk}{0.41}
\newcommand{\LowMlRfQwk}{0.26}
\newcommand{\HighMlRfQwk}{0.38}
\newcommand{\ObsMlRfQwk}{2.65}
\newcommand{\ExpMlRfQwk}{3.36}
\newcommand{\QwkMlRfQwk}{0.211}
\newcommand{\NamaeMlRfQwk}{0.326}
\newcommand{\NmmaeMlRfQwk}{0.399}
\newcommand{\PredVarMlRfAmae}{1.27}
\newcommand{\EndShareMlRfAmae}{0.28}
\newcommand{\LowMlRfAmae}{0.14}
\newcommand{\HighMlRfAmae}{0.32}
\newcommand{\ObsMlRfAmae}{2.04}
\newcommand{\ExpMlRfAmae}{2.55}
\newcommand{\QwkMlRfAmae}{0.202}
\newcommand{\NamaeMlRfAmae}{0.321}
\newcommand{\NmmaeMlRfAmae}{0.527}
\newcommand{\PredVarWrPoQwk}{0.37}
\newcommand{\EndShareWrPoQwk}{0.00}
\newcommand{\LowWrPoQwk}{0.00}
\newcommand{\HighWrPoQwk}{0.00}
\newcommand{\ObsWrPoQwk}{0.51}
\newcommand{\ExpWrPoQwk}{1.03}
\newcommand{\QwkWrPoQwk}{0.502}
\newcommand{\NamaeWrPoQwk}{0.212}
\newcommand{\NmmaeWrPoQwk}{0.410}
\newcommand{\PredVarWrPoAmae}{1.92}
\newcommand{\EndShareWrPoAmae}{0.16}
\newcommand{\LowWrPoAmae}{0.45}
\newcommand{\HighWrPoAmae}{0.61}
\newcommand{\ObsWrPoAmae}{1.37}
\newcommand{\ExpWrPoAmae}{2.58}
\newcommand{\QwkWrPoAmae}{0.469}
\newcommand{\NamaeWrPoAmae}{0.154}
\newcommand{\NmmaeWrPoAmae}{0.193}
\newcommand{\LabVarWr}{0.65}
\newcommand{\LabEndWr}{0.02}
\newcommand{\PredVarWrRfQwk}{0.80}
\newcommand{\EndShareWrRfQwk}{0.02}
\newcommand{\LowWrRfQwk}{0.00}
\newcommand{\HighWrRfQwk}{0.06}
\newcommand{\ObsWrRfQwk}{0.67}
\newcommand{\ExpWrRfQwk}{1.45}
\newcommand{\QwkWrRfQwk}{0.541}
\newcommand{\NamaeWrRfQwk}{0.173}
\newcommand{\NmmaeWrRfQwk}{0.320}
\newcommand{\PredVarWrRfAmae}{1.36}
\newcommand{\EndShareWrRfAmae}{0.07}
\newcommand{\LowWrRfAmae}{0.05}
\newcommand{\HighWrRfAmae}{0.22}
\newcommand{\ObsWrRfAmae}{1.04}
\newcommand{\ExpWrRfAmae}{2.01}
\newcommand{\QwkWrRfAmae}{0.484}
\newcommand{\NamaeWrRfAmae}{0.168}
\newcommand{\NmmaeWrRfAmae}{0.270}
\newcommand{\PredVarScPoQwk}{2.73}
\newcommand{\EndShareScPoQwk}{0.14}
\newcommand{\LowScPoQwk}{0.56}
\newcommand{\HighScPoQwk}{0.70}
\newcommand{\ObsScPoQwk}{1.32}
\newcommand{\ExpScPoQwk}{4.83}
\newcommand{\QwkScPoQwk}{0.726}
\newcommand{\NamaeScPoQwk}{0.124}
\newcommand{\NmmaeScPoQwk}{0.155}
\newcommand{\PredVarScPoAmae}{2.73}
\newcommand{\EndShareScPoAmae}{0.13}
\newcommand{\LowScPoAmae}{0.56}
\newcommand{\HighScPoAmae}{0.70}
\newcommand{\ObsScPoAmae}{1.32}
\newcommand{\ExpScPoAmae}{4.83}
\newcommand{\QwkScPoAmae}{0.726}
\newcommand{\NamaeScPoAmae}{0.124}
\newcommand{\NmmaeScPoAmae}{0.155}
\newcommand{\LabVarSc}{2.10}
\newcommand{\LabEndSc}{0.06}
\newcommand{\PredVarScRfQwk}{2.39}
\newcommand{\EndShareScRfQwk}{0.07}
\newcommand{\LowScRfQwk}{0.46}
\newcommand{\HighScRfQwk}{0.14}
\newcommand{\ObsScRfQwk}{1.19}
\newcommand{\ExpScRfQwk}{4.48}
\newcommand{\QwkScRfQwk}{0.735}
\newcommand{\NamaeScRfQwk}{0.131}
\newcommand{\NmmaeScRfQwk}{0.151}
\newcommand{\PredVarScRfAmae}{2.87}
\newcommand{\EndShareScRfAmae}{0.11}
\newcommand{\LowScRfAmae}{0.54}
\newcommand{\HighScRfAmae}{0.44}
\newcommand{\ObsScRfAmae}{1.42}
\newcommand{\ExpScRfAmae}{4.97}
\newcommand{\QwkScRfAmae}{0.714}
\newcommand{\NamaeScRfAmae}{0.132}
\newcommand{\NmmaeScRfAmae}{0.174}
\newcommand{\ConstNamaeMl}{0.300}
\newcommand{\BestCandNamaeMl}{0.320}
\newcommand{\ConstNmaeMl}{0.253}
\newcommand{\ConstNamaeRe}{0.333}
\newcommand{\BestCandNamaeRe}{0.407}
\newcommand{\ConstNmaeRe}{0.319}
\newcommand{\ConstNamaeSu}{0.300}
\newcommand{\BestCandNamaeSu}{0.336}
\newcommand{\ConstNmaeSu}{0.301}
\newcommand{\QwkRangeLow}{0.07}
\newcommand{\QwkRangeHigh}{0.21}
\newcommand{\SimAmaeMin}{\ensuremath{-}0.077}
\newcommand{\SimAmaeMax}{\ensuremath{+}0.072}
\newcommand{\SimAmaeAbsSmall}{9}
\newcommand{\SimCells}{12}


\title{Online Resource 1: Supplementary Information\\[4pt]\large Which metric should choose the model? Selection criteria, class weighting and endpoint recall in tabular ordinal classification}
\author{Muhammed Yusuf Batur \and Mehmet G\"{o}kt\"{u}rk}
\date{}

\begin{document}
\maketitle

\begin{center}\small
Journal: Computational Statistics\\
Corresponding author: Muhammed Yusuf Batur, baturmuhammedyusuf@gmail.com\\
Department of Computer Sciences, Faculty of Engineering and Natural Sciences, Istanbul Rumeli University, Istanbul, T\"{u}rkiye\\
Department of Computer Engineering, Graduate School, Gebze Technical University, Kocaeli, T\"{u}rkiye
\end{center}

\section*{S1 Data provenance, features and dependence}

All tasks are secondary analyses of publicly documented data. No raw records, identifiers or row-level predictions are distributed with this article.

\paragraph{Wine (red, white).} UCI Wine Quality (DOI 10.24432/C56S3T, CC BY 4.0; Cortez et al.\ 2009a). Response: sensory quality score (red 3--8, white 3--9). Predictors: the 11 physicochemical variables. Groups: identical predictor vectors (1,359 groups for 1,599 red wines; 3,961 for 4,898 white wines). Latent production batches are unknown. The two colours form one source family. Endpoint support is very small (red 10 and 18 cases; white 20 and 5).

\paragraph{SkillCraft.} UCI SkillCraft1 Master Table (DOI 10.24432/C5161N, CC BY 4.0; Thompson et al.\ 2013, 2014). Response: league index 1--7 (Bronze to Grandmaster). League 8 (55 professional games, collected from gaming websites rather than from survey participants) was excluded on provenance grounds, decided before the restricted task was fitted; the original eight-league task is reported below as superseded. Predictors: 18 numeric telemetry and self-report variables (GameID excluded). Groups: identical predictor vectors (3,339 groups for 3,340 records). The source studies report one replay per survey participant; this cannot be verified record by record, so independence of persons is an assumption.

\paragraph{Restaurant.} UCI Restaurant \& consumer data (DOI 10.24432/C5DP41, CC BY 4.0). Response: overall rating 0--2. Predictors (all categorical): restaurant alcohol, dress code, accessibility, price, ambience, franchise, area, other services; user budget, dress preference, ambience, transport. Excluded: food and service ratings, identifiers, names, locations, demographics, religion, marital status and weight. Groups: users (138). The UCI release (1,161 ratings, 138 users) differs from the counts reported in the original paper; we do not claim that they are the same version.

\paragraph{Sushi.} Sushi preference data, 2016 corrected release (Kamishima 2003), research use permitted, redistribution not permitted. Response: score 0 (most disliked) to 4 (most preferred); unrated entries ($-1$) removed. 500 of 5,000 users were selected by the lowest SHA-256 hash of a fixed salt and the user identifier, independently of outcomes; all their observed scores were retained (5,000 scores). Predictors: user age band, item style, major and minor group. Groups: users.

\paragraph{MovieLens-100k.} GroupLens MovieLens-100k (Harper and Konstan 2015), research use with attribution, no redistribution and no endorsement. Response: rating 1--5 mapped to ranks 0--4. 400 of 943 users were selected by the same hash rule; all their ratings were retained (43,699 ratings). Predictors: user age, movie release year and the 19 genre indicators. Excluded: identifiers, gender, occupation, ZIP code, title, URL, timestamp and rating-derived summaries. Groups: users.

\paragraph{TAE (control).} UCI Teaching Assistant Evaluation (DOI 10.24432/C55P6M, CC BY 4.0). Response: low/medium/high, created by the source by binning scores; therefore a control task, not a primary outcome. Predictors: native English speaker, course, summer or regular semester, class size. Groups: course instructor (25 groups).

\paragraph{CTG (stage 3).} UCI Cardiotocography (DOI 10.24432/C51S4N, CC BY 4.0; Ayres-de-Campos et al.\ 2000). Response: fetal state NSP (normal $<$ suspect $<$ pathologic), assigned by consensus of three expert obstetricians. Predictors: the 21 automatically measured per-unit-time signal features; the expert baseline value and the morphological class labels are not used. Groups: examination file (352 files for 2,126 segments); the file is not a verified patient identifier.

\paragraph{HCV (stage 3).} UCI HCV data (DOI 10.24432/C5D612, CC BY 4.0; Hoffmann et al.\ 2018). Response: blood donor $<$ hepatitis $<$ fibrosis $<$ cirrhosis, the progression order given by the source. The seven ``suspect blood donors'' were excluded before fitting because they are not a documented level of the ordered scale; they were not merged with the donors. Predictors: age, sex and ten laboratory values. Groups: patient (one record per patient). The data combine donors and patients by design (case--control sampling).

\paragraph{Cirrhosis (stage 3).} UCI Cirrhosis Patient Survival Prediction (DOI 10.24432/C5R02G, CC BY 4.0; Dickson et al.\ 1989b), Mayo Clinic primary biliary cirrhosis study (312 participants of a randomized trial and 106 further patients who consented to basic measurements and follow-up). Response: histologic stage 1--4; six patients without a stage were excluded. Predictors: 15 baseline clinical and laboratory variables; follow-up time, status and the randomized drug assignment are not used. Groups: patient.

\paragraph{SUPPORT2 (stage 3).} SUPPORT study data (Knaus et al.\ 1995) as distributed by the Vanderbilt Department of Biostatistics (hbiostat.org/data), whose terms grant permission for use; the UCI copy (dataset 880) states no licence, and no redistribution right is claimed. Response: the five-level functional disability scale sfdm2 in its documented order. 5,000 of the 7,705 labelled patients were selected by SHA-256 hash of the identifier. Predictors: 29 baseline variables; the APS/SPS scores, model-based survival and prognosis estimates, do-not-resuscitate variables, follow-up and cost variables are not used. Groups: patient.

\paragraph{Control tasks added in stage 3.} Diabetes130 (UCI DOI 10.24432/C5230J, CC BY 4.0; Strack et al.\ 2014): readmission NO $<$ $>$30 days $<$ $<$30 days, a source binning of the time to readmission; encounters ending in death or hospice excluded; at most 5,000 encounters from 3,471 patients selected by patient hash; groups: patient. Student (UCI DOI 10.24432/C5TG7T, CC BY 4.0): final Portuguese grade binned by the source's five-level scheme (0--9, 10--11, 12--13, 14--15, 16--20); intermediate grades not used. NPHA (ICPSR DOI 10.3886/ICPSR37305.v1; UCI dataset 936, CC BY 4.0): number of doctors visited, binned as 0--1, 2--3 and 4 or more; all 14 survey items categorical.

\paragraph{Sources excluded in stage 3 before any fit.} Maternal Health Risk (the label is defined by the source as a predicted risk level derived from the predictors), CDC Diabetes Health Indicators (no licence stated; undocumented derivation of the table), Higher Education Students Performance (lowest category has 8 students, below the 20-group minimum) and Heart Disease (order and meaning of the severity levels 1--4 undocumented; sites cannot be pooled). No excluded source was replaced.

\paragraph{Item overlap.} In the three rating tasks, items recur across users: roughly 99--100\% of test ratings concern items that also occur in the training partition. Results therefore refer to new users on the observed item pool.

\paragraph{Stage 5: Drug consumption (alcohol).} UCI 373 Drug Consumption (Quantified), DOI 10.24432/C5TC7S, CC BY 4.0 (Fehrman et al.\ 2017). 1,885 respondents, one row each; group = respondent ID. Target: alcohol use CL0 never used $<$ CL1 over a decade ago $<$ CL2 in the last decade $<$ CL3 last year $<$ CL4 last month $<$ CL5 last week $<$ CL6 last day. Alcohol was fixed before any outcome was seen as the first substance in the source's variable list. Predictors: the 12 quantified attributes (age, gender, education, country, ethnicity, five NEO-FFI-R scores, impulsivity, sensation seeking). The other 17 substances and the fictitious control drug are excluded.
\paragraph{Stage 5: Entrance examination.} UCI 582 Student Performance on an Entrance Examination, DOI 10.24432/C58D0H, CC BY 4.0. 666 candidates; target: performance Average $<$ Good $<$ Very good $<$ Excellent; predictors: the 11 categorical attributes; groups: identical predictor vectors (576 groups).
\paragraph{Stage 5: NHANES self-rated health.} NHANES 2017--2018 public-use files DEMO\_J, HSQ\_J, BMX\_J, BPX\_J, TCHOL\_J, HDL\_J and GHB\_J, merged by SEQN, used under the NCHS Data User Agreement (statistical analysis only, no re-identification, no linkage with identifiable data). Adults aged 20 or more with HSD010 between 1 (excellent) and 5 (poor); refusals and unknowns are excluded (4,890 adults). Predictors: age, sex, race/ethnicity, education, income-to-poverty ratio, BMI, waist circumference, mean systolic and diastolic blood pressure, total cholesterol, HDL cholesterol and HbA1c.
\paragraph{Stage 5: NHIS self-rated health.} NHIS 2019 sample adult file, same agreement. PHSTAT\_A from 1 (excellent) to 5 (poor); codes 7--9 are excluded; one sample adult per household (group = HHX); outcome-blind hash subsample of 5,000 adults. Predictors: age, poverty ratio, sex, Hispanic origin and race, education, region, urbanicity, marital status, BMI category, smoking status and twelve ever-diagnosed conditions. Survey weights are not used.
\paragraph{Stage 5: blocked sources.} Three sources failed the rules before download: GSS 2018 (no stated terms of use), OpenML eucalyptus (rating determined from measured predictors) and UCI thyroid (nominal coding). Six failed the endpoint-support rule after download: OpenML SWD (17 groups in the lowest class), LEV (14 in the highest), ERA (4) and ESL (2), UCI 467 (8 students in the highest class) and UCI 397 (11 reviews with score 1).

\section*{S2 Protocol and verification}

Splits: stratified group $K$-fold, three outer and two inner folds, two repeats (seeds 20261120/20261121 for the first stage, 20261124/20261125 for the second and third stages); for each seed, the first support-valid assignment among offsets 0--99 was used (all categories present in every partition, no group overlap). Candidates: proportional odds with $C\in\{0.1,1\}$, random forest with depth $\in\{4,\text{unlimited}\}$, each with and without balanced weights. Endpoint-recall floor 0.2 with AMAE fallback; tie tolerance $10^{-12}$. Protocols, configurations, data and split hashes were frozen locally (SHA-256) before the fits of each stage; there was no external registration. An independent re-implementation of the metrics reproduced the stored values (\NRecompute{} comparisons, largest absolute difference \RecomputeMaxDiff{}).

\paragraph{Stage 4 (pre-specified robustness).} The protocol (file \texttt{robustness\_m2/PROTOCOL.md}), the configuration, the runner code and the hashes of the reused split assignments were frozen with SHA-256 hashes and a UTC time stamp before any stage-4 fit; there was no external registration. Part A refits the 14 tasks on the frozen splits of stages 1--3 with 24 candidates: proportional odds with $C\in\{0.01,0.1,1,10\}$ (at most 1000 iterations); random forest with 100 trees, minimum leaf size 2 and depth $\in\{3,6,12,\text{unlimited}\}$; histogram gradient boosting (scikit-learn \texttt{HistGradientBoostingClassifier}) with (learning rate, maximum depth) $\in\{(0.05,3),(0.05,\text{unlimited}),(0.1,3),(0.1,\text{unlimited})\}$, 200 iterations, at most 31 leaves and no early stopping; each with and without balanced weights (\FitsStageFour{} fits). Part B simulates $K=5$ categories from the latent Gaussian model with lowest-category share $\pi_0\in\{0.02,0.05,0.10,0.20,0.30,0.40\}$, highest-category share 0.20, the three middle categories sharing the remainder, $\lambda\in\{0.5,1,2\}$, ten replications (seeds 20261301--20261310) and the original eight candidates (\FitsSimFour{} fits). No candidate, task or criterion was dropped and no fit failed.

\paragraph{Stage 5 (pre-registered validation).} The protocol (\texttt{validation\_m5/PROTOCOL.md}), configuration, code, file hashes, eligibility audit and split assignments were frozen and pushed to the public repository (commit \texttt{66e426e}) before any stage-5 fit. Candidates: the 24 stage-4 candidates and a Frank--Hall family ($K-1$ binary histogram gradient-boosting models for $P(y>k)$ with the four stage-4 settings, class probabilities by differences, clipped at zero and renormalized; balanced weights $n/(K n_y)$ passed to every binary model). Seeds 20261124/20261125 for the new tasks; the frozen splits of stages 1--3 for the others (\FitsStageFive{} fits). Median decoding predicts the smallest $k$ with cumulative probability at least 0.5. Floor with max--min fallback: the largest inner QWK among candidates with both inner endpoint recalls at least 0.2; if none, the largest smaller inner endpoint recall, then the larger inner QWK, then candidate order. The 24 stage-4 candidates were reproduced exactly on the 14 earlier tasks.

\section*{S3 Tables}

{\footnotesize
\setlength{\tabcolsep}{3.5pt}
\begin{longtable}{@{}lll*{9}{r}@{}}
\caption{Outer-test metrics of the policy selected by each criterion (mean of two repeats). Error metrics are normalized by $K-1$. Bal.: fraction of decisions selecting balanced weights.}\label{tab:s-policies}\\
\toprule
Task & Fam. & Criterion & QWK & MAE & AMAE & MMAE & Macro-F1 & RPS & $R_0$ & $R_{K-1}$ & Bal.\\
\midrule\endfirsthead
\toprule
Task & Fam. & Criterion & QWK & MAE & AMAE & MMAE & Macro-F1 & RPS & $R_0$ & $R_{K-1}$ & Bal.\\
\midrule\endhead
\bottomrule\endfoot
Wine (red) & PO & QWK & 0.502 & 0.089 & 0.212 & 0.410 & 0.269 & 0.064 & 0.00 & 0.00 & 0.00 \\
 &  & MAE & 0.499 & 0.088 & 0.214 & 0.420 & 0.268 & 0.064 & 0.00 & 0.00 & 0.00 \\
 &  & AMAE & 0.469 & 0.177 & 0.154 & 0.193 & 0.237 & 0.111 & 0.45 & 0.61 & 1.00 \\
 &  & MMAE & 0.470 & 0.177 & 0.153 & 0.192 & 0.237 & 0.111 & 0.45 & 0.61 & 1.00 \\
 &  & Macro-F1 & 0.502 & 0.089 & 0.212 & 0.410 & 0.269 & 0.064 & 0.00 & 0.00 & 0.00 \\
 &  & RPS & 0.499 & 0.089 & 0.213 & 0.410 & 0.267 & 0.064 & 0.00 & 0.00 & 0.00 \\
 &  & QWK-floor & 0.471 & 0.177 & 0.153 & 0.192 & 0.237 & 0.110 & 0.45 & 0.61 & 1.00 \\
\addlinespace
 & RF & QWK & 0.541 & 0.104 & 0.173 & 0.320 & 0.312 & 0.068 & 0.00 & 0.06 & 1.00 \\
 &  & MAE & 0.492 & 0.088 & 0.226 & 0.450 & 0.270 & 0.062 & 0.00 & 0.00 & 0.00 \\
 &  & AMAE & 0.484 & 0.141 & 0.168 & 0.270 & 0.269 & 0.082 & 0.05 & 0.22 & 1.00 \\
 &  & MMAE & 0.484 & 0.141 & 0.168 & 0.270 & 0.269 & 0.082 & 0.05 & 0.22 & 1.00 \\
 &  & Macro-F1 & 0.541 & 0.104 & 0.173 & 0.320 & 0.312 & 0.068 & 0.00 & 0.06 & 1.00 \\
 &  & RPS & 0.492 & 0.088 & 0.226 & 0.450 & 0.270 & 0.062 & 0.00 & 0.00 & 0.00 \\
 &  & QWK-floor & 0.484 & 0.141 & 0.168 & 0.270 & 0.269 & 0.082 & 0.05 & 0.22 & 1.00 \\
\addlinespace
Wine (white) & PO & QWK & 0.401 & 0.089 & 0.226 & 0.438 & 0.213 & 0.063 & 0.00 & 0.00 & 0.00 \\
 &  & MAE & 0.401 & 0.089 & 0.226 & 0.438 & 0.213 & 0.063 & 0.00 & 0.00 & 0.00 \\
 &  & AMAE & 0.378 & 0.130 & 0.236 & 0.458 & 0.209 & 0.079 & 0.10 & 0.10 & 0.33 \\
 &  & MMAE & 0.375 & 0.110 & 0.242 & 0.438 & 0.208 & 0.072 & 0.07 & 0.00 & 0.17 \\
 &  & Macro-F1 & 0.401 & 0.089 & 0.226 & 0.438 & 0.213 & 0.063 & 0.00 & 0.00 & 0.00 \\
 &  & RPS & 0.401 & 0.089 & 0.226 & 0.438 & 0.213 & 0.063 & 0.00 & 0.00 & 0.00 \\
 &  & QWK-floor & 0.368 & 0.187 & 0.234 & 0.450 & 0.177 & 0.101 & 0.28 & 0.50 & 0.83 \\
\addlinespace
 & RF & QWK & 0.554 & 0.095 & 0.193 & 0.400 & 0.299 & 0.063 & 0.07 & 0.00 & 1.00 \\
 &  & MAE & 0.464 & 0.084 & 0.227 & 0.450 & 0.241 & 0.059 & 0.00 & 0.00 & 0.00 \\
 &  & AMAE & 0.554 & 0.095 & 0.193 & 0.400 & 0.299 & 0.063 & 0.07 & 0.00 & 1.00 \\
 &  & MMAE & 0.454 & 0.127 & 0.208 & 0.396 & 0.260 & 0.076 & 0.10 & 0.10 & 1.00 \\
 &  & Macro-F1 & 0.554 & 0.095 & 0.193 & 0.400 & 0.299 & 0.063 & 0.07 & 0.00 & 1.00 \\
 &  & RPS & 0.487 & 0.082 & 0.225 & 0.446 & 0.253 & 0.059 & 0.00 & 0.00 & 0.00 \\
 &  & QWK-floor & 0.554 & 0.095 & 0.193 & 0.400 & 0.299 & 0.063 & 0.07 & 0.00 & 1.00 \\
\addlinespace
SkillCraft & PO & QWK & 0.726 & 0.139 & 0.124 & 0.155 & 0.357 & 0.094 & 0.56 & 0.70 & 1.00 \\
 &  & MAE & 0.706 & 0.124 & 0.146 & 0.212 & 0.329 & 0.086 & 0.24 & 0.03 & 0.00 \\
 &  & AMAE & 0.726 & 0.139 & 0.124 & 0.155 & 0.358 & 0.094 & 0.56 & 0.70 & 1.00 \\
 &  & MMAE & 0.725 & 0.139 & 0.124 & 0.155 & 0.358 & 0.094 & 0.56 & 0.70 & 1.00 \\
 &  & Macro-F1 & 0.725 & 0.140 & 0.124 & 0.155 & 0.355 & 0.094 & 0.55 & 0.70 & 1.00 \\
 &  & RPS & 0.706 & 0.124 & 0.146 & 0.212 & 0.330 & 0.086 & 0.24 & 0.03 & 0.00 \\
 &  & QWK-floor & 0.726 & 0.139 & 0.124 & 0.155 & 0.357 & 0.094 & 0.56 & 0.70 & 1.00 \\
\addlinespace
 & RF & QWK & 0.735 & 0.130 & 0.131 & 0.151 & 0.364 & 0.084 & 0.46 & 0.14 & 1.00 \\
 &  & MAE & 0.710 & 0.124 & 0.145 & 0.204 & 0.333 & 0.084 & 0.28 & 0.00 & 0.00 \\
 &  & AMAE & 0.714 & 0.144 & 0.132 & 0.174 & 0.345 & 0.090 & 0.54 & 0.44 & 1.00 \\
 &  & MMAE & 0.735 & 0.130 & 0.131 & 0.151 & 0.364 & 0.084 & 0.46 & 0.14 & 1.00 \\
 &  & Macro-F1 & 0.731 & 0.133 & 0.131 & 0.155 & 0.365 & 0.085 & 0.49 & 0.23 & 1.00 \\
 &  & RPS & 0.716 & 0.123 & 0.143 & 0.191 & 0.343 & 0.082 & 0.31 & 0.00 & 0.00 \\
 &  & QWK-floor & 0.712 & 0.147 & 0.131 & 0.179 & 0.339 & 0.092 & 0.56 & 0.54 & 1.00 \\
\addlinespace
Restaurant & PO & QWK & 0.075 & 0.380 & 0.424 & 0.669 & 0.358 & 0.218 & 0.17 & 0.58 & 0.33 \\
 &  & MAE & 0.080 & 0.343 & 0.407 & 0.764 & 0.350 & 0.211 & 0.04 & 0.66 & 0.00 \\
 &  & AMAE & 0.084 & 0.363 & 0.414 & 0.705 & 0.358 & 0.215 & 0.12 & 0.62 & 0.17 \\
 &  & MMAE & 0.052 & 0.442 & 0.451 & 0.507 & 0.336 & 0.226 & 0.40 & 0.45 & 1.00 \\
 &  & Macro-F1 & 0.072 & 0.399 & 0.429 & 0.600 & 0.360 & 0.221 & 0.27 & 0.53 & 0.50 \\
 &  & RPS & 0.080 & 0.343 & 0.407 & 0.764 & 0.350 & 0.211 & 0.04 & 0.66 & 0.00 \\
 &  & QWK-floor & 0.046 & 0.444 & 0.453 & 0.512 & 0.335 & 0.225 & 0.39 & 0.45 & 1.00 \\
\addlinespace
 & RF & QWK & 0.074 & 0.371 & 0.420 & 0.695 & 0.360 & 0.214 & 0.13 & 0.63 & 0.33 \\
 &  & MAE & 0.081 & 0.348 & 0.415 & 0.783 & 0.347 & 0.209 & 0.04 & 0.73 & 0.00 \\
 &  & AMAE & 0.064 & 0.376 & 0.422 & 0.679 & 0.363 & 0.218 & 0.14 & 0.60 & 0.33 \\
 &  & MMAE & 0.054 & 0.416 & 0.427 & 0.520 & 0.361 & 0.221 & 0.32 & 0.44 & 1.00 \\
 &  & Macro-F1 & 0.072 & 0.381 & 0.420 & 0.647 & 0.364 & 0.220 & 0.17 & 0.59 & 0.50 \\
 &  & RPS & 0.094 & 0.339 & 0.412 & 0.809 & 0.343 & 0.207 & 0.01 & 0.77 & 0.00 \\
 &  & QWK-floor & 0.041 & 0.417 & 0.432 & 0.533 & 0.359 & 0.225 & 0.31 & 0.45 & 1.00 \\
\addlinespace
Sushi & PO & QWK & 0.176 & 0.382 & 0.377 & 0.398 & 0.211 & 0.194 & 0.45 & 0.45 & 1.00 \\
 &  & MAE & 0.116 & 0.287 & 0.413 & 0.801 & 0.167 & 0.164 & 0.00 & 0.81 & 0.00 \\
 &  & AMAE & 0.176 & 0.382 & 0.377 & 0.398 & 0.211 & 0.194 & 0.45 & 0.45 & 1.00 \\
 &  & MMAE & 0.176 & 0.381 & 0.377 & 0.397 & 0.212 & 0.194 & 0.45 & 0.45 & 1.00 \\
 &  & Macro-F1 & 0.176 & 0.384 & 0.378 & 0.400 & 0.209 & 0.194 & 0.46 & 0.45 & 1.00 \\
 &  & RPS & 0.105 & 0.291 & 0.422 & 0.819 & 0.162 & 0.164 & 0.00 & 0.83 & 0.00 \\
 &  & QWK-floor & 0.176 & 0.382 & 0.377 & 0.398 & 0.211 & 0.194 & 0.45 & 0.45 & 1.00 \\
\addlinespace
 & RF & QWK & 0.163 & 0.354 & 0.344 & 0.420 & 0.239 & 0.196 & 0.33 & 0.32 & 1.00 \\
 &  & MAE & 0.092 & 0.288 & 0.415 & 0.819 & 0.198 & 0.166 & 0.02 & 0.75 & 0.00 \\
 &  & AMAE & 0.170 & 0.350 & 0.339 & 0.420 & 0.242 & 0.196 & 0.32 & 0.31 & 1.00 \\
 &  & MMAE & 0.162 & 0.354 & 0.346 & 0.416 & 0.237 & 0.197 & 0.32 & 0.32 & 1.00 \\
 &  & Macro-F1 & 0.170 & 0.350 & 0.339 & 0.420 & 0.242 & 0.196 & 0.32 & 0.31 & 1.00 \\
 &  & RPS & 0.035 & 0.307 & 0.470 & 0.942 & 0.148 & 0.165 & 0.00 & 0.93 & 0.00 \\
 &  & QWK-floor & 0.163 & 0.354 & 0.344 & 0.420 & 0.239 & 0.196 & 0.33 & 0.32 & 1.00 \\
\addlinespace
MovieLens-100k & PO & QWK & 0.196 & 0.367 & 0.354 & 0.388 & 0.207 & 0.172 & 0.52 & 0.40 & 1.00 \\
 &  & MAE & 0.137 & 0.218 & 0.323 & 0.669 & 0.182 & 0.149 & 0.00 & 0.11 & 0.00 \\
 &  & AMAE & 0.137 & 0.218 & 0.323 & 0.669 & 0.182 & 0.149 & 0.00 & 0.11 & 0.00 \\
 &  & MMAE & 0.196 & 0.367 & 0.354 & 0.388 & 0.207 & 0.172 & 0.52 & 0.40 & 1.00 \\
 &  & Macro-F1 & 0.196 & 0.367 & 0.354 & 0.388 & 0.207 & 0.172 & 0.52 & 0.40 & 1.00 \\
 &  & RPS & 0.137 & 0.218 & 0.323 & 0.669 & 0.182 & 0.149 & 0.00 & 0.11 & 0.00 \\
 &  & QWK-floor & 0.196 & 0.367 & 0.354 & 0.388 & 0.207 & 0.172 & 0.52 & 0.40 & 1.00 \\
\addlinespace
 & RF & QWK & 0.211 & 0.311 & 0.326 & 0.399 & 0.249 & 0.171 & 0.26 & 0.38 & 1.00 \\
 &  & MAE & 0.016 & 0.223 & 0.348 & 0.745 & 0.114 & 0.149 & 0.00 & 0.02 & 0.00 \\
 &  & AMAE & 0.202 & 0.265 & 0.321 & 0.527 & 0.259 & 0.161 & 0.14 & 0.32 & 0.33 \\
 &  & MMAE & 0.192 & 0.354 & 0.351 & 0.369 & 0.211 & 0.171 & 0.32 & 0.47 & 1.00 \\
 &  & Macro-F1 & 0.208 & 0.288 & 0.322 & 0.456 & 0.258 & 0.166 & 0.21 & 0.35 & 0.67 \\
 &  & RPS & 0.016 & 0.223 & 0.348 & 0.745 & 0.114 & 0.149 & 0.00 & 0.02 & 0.00 \\
 &  & QWK-floor & 0.194 & 0.337 & 0.343 & 0.390 & 0.227 & 0.170 & 0.29 & 0.44 & 1.00 \\
\addlinespace
CTG & PO & QWK & 0.736 & 0.086 & 0.175 & 0.250 & 0.686 & 0.060 & 0.92 & 0.57 & 0.17 \\
 &  & MAE & 0.736 & 0.081 & 0.193 & 0.284 & 0.670 & 0.056 & 0.95 & 0.52 & 0.00 \\
 &  & AMAE & 0.709 & 0.110 & 0.123 & 0.142 & 0.699 & 0.078 & 0.80 & 0.76 & 1.00 \\
 &  & MMAE & 0.709 & 0.110 & 0.123 & 0.142 & 0.699 & 0.078 & 0.80 & 0.76 & 1.00 \\
 &  & Macro-F1 & 0.697 & 0.110 & 0.138 & 0.181 & 0.690 & 0.078 & 0.82 & 0.70 & 0.83 \\
 &  & RPS & 0.741 & 0.080 & 0.190 & 0.281 & 0.675 & 0.056 & 0.95 & 0.54 & 0.00 \\
 &  & QWK-floor & 0.736 & 0.086 & 0.175 & 0.250 & 0.686 & 0.060 & 0.92 & 0.57 & 0.17 \\
\addlinespace
 & RF & QWK & 0.813 & 0.053 & 0.109 & 0.157 & 0.831 & 0.040 & 0.96 & 0.78 & 0.00 \\
 &  & MAE & 0.813 & 0.053 & 0.109 & 0.157 & 0.831 & 0.040 & 0.96 & 0.78 & 0.00 \\
 &  & AMAE & 0.772 & 0.067 & 0.083 & 0.102 & 0.823 & 0.053 & 0.90 & 0.86 & 1.00 \\
 &  & MMAE & 0.766 & 0.069 & 0.084 & 0.101 & 0.818 & 0.057 & 0.90 & 0.85 & 1.00 \\
 &  & Macro-F1 & 0.791 & 0.059 & 0.094 & 0.125 & 0.831 & 0.044 & 0.94 & 0.83 & 0.50 \\
 &  & RPS & 0.813 & 0.053 & 0.109 & 0.157 & 0.831 & 0.040 & 0.96 & 0.78 & 0.00 \\
 &  & QWK-floor & 0.813 & 0.053 & 0.109 & 0.157 & 0.831 & 0.040 & 0.96 & 0.78 & 0.00 \\
\addlinespace
HCV & PO & QWK & 0.745 & 0.046 & 0.279 & 0.468 & 0.477 & 0.033 & 1.00 & 0.70 & 0.00 \\
 &  & MAE & 0.731 & 0.047 & 0.285 & 0.468 & 0.478 & 0.033 & 1.00 & 0.70 & 0.00 \\
 &  & AMAE & 0.639 & 0.130 & 0.166 & 0.257 & 0.494 & 0.111 & 0.66 & 0.75 & 1.00 \\
 &  & MMAE & 0.639 & 0.130 & 0.166 & 0.257 & 0.494 & 0.111 & 0.66 & 0.75 & 1.00 \\
 &  & Macro-F1 & 0.684 & 0.081 & 0.235 & 0.365 & 0.504 & 0.065 & 0.86 & 0.75 & 0.33 \\
 &  & RPS & 0.695 & 0.051 & 0.313 & 0.556 & 0.433 & 0.034 & 1.00 & 0.68 & 0.00 \\
 &  & QWK-floor & 0.745 & 0.046 & 0.279 & 0.468 & 0.477 & 0.033 & 1.00 & 0.70 & 0.00 \\
\addlinespace
 & RF & QWK & 0.841 & 0.035 & 0.151 & 0.286 & 0.635 & 0.034 & 0.98 & 0.83 & 0.83 \\
 &  & MAE & 0.817 & 0.038 & 0.176 & 0.325 & 0.620 & 0.032 & 0.98 & 0.82 & 0.50 \\
 &  & AMAE & 0.785 & 0.049 & 0.143 & 0.248 & 0.604 & 0.051 & 0.94 & 0.77 & 1.00 \\
 &  & MMAE & 0.814 & 0.043 & 0.135 & 0.232 & 0.624 & 0.049 & 0.95 & 0.82 & 1.00 \\
 &  & Macro-F1 & 0.822 & 0.040 & 0.140 & 0.254 & 0.624 & 0.042 & 0.96 & 0.82 & 1.00 \\
 &  & RPS & 0.789 & 0.038 & 0.224 & 0.365 & 0.565 & 0.024 & 1.00 & 0.73 & 0.00 \\
 &  & QWK-floor & 0.841 & 0.035 & 0.151 & 0.286 & 0.635 & 0.034 & 0.98 & 0.83 & 0.83 \\
\addlinespace
Cirrhosis & PO & QWK & 0.382 & 0.265 & 0.287 & 0.357 & 0.342 & 0.158 & 0.33 & 0.59 & 0.83 \\
 &  & MAE & 0.377 & 0.202 & 0.295 & 0.587 & 0.348 & 0.135 & 0.00 & 0.61 & 0.00 \\
 &  & AMAE & 0.380 & 0.282 & 0.294 & 0.333 & 0.328 & 0.166 & 0.43 & 0.58 & 1.00 \\
 &  & MMAE & 0.380 & 0.283 & 0.295 & 0.341 & 0.320 & 0.166 & 0.38 & 0.58 & 1.00 \\
 &  & Macro-F1 & 0.383 & 0.274 & 0.300 & 0.381 & 0.318 & 0.160 & 0.36 & 0.59 & 0.83 \\
 &  & RPS & 0.380 & 0.199 & 0.293 & 0.587 & 0.353 & 0.134 & 0.00 & 0.61 & 0.00 \\
 &  & QWK-floor & 0.383 & 0.281 & 0.292 & 0.331 & 0.325 & 0.163 & 0.40 & 0.58 & 1.00 \\
\addlinespace
 & RF & QWK & 0.426 & 0.231 & 0.264 & 0.357 & 0.399 & 0.144 & 0.33 & 0.69 & 1.00 \\
 &  & MAE & 0.351 & 0.198 & 0.301 & 0.611 & 0.364 & 0.135 & 0.05 & 0.67 & 0.17 \\
 &  & AMAE & 0.431 & 0.229 & 0.257 & 0.333 & 0.409 & 0.145 & 0.38 & 0.70 & 1.00 \\
 &  & MMAE & 0.440 & 0.230 & 0.251 & 0.318 & 0.418 & 0.146 & 0.45 & 0.72 & 1.00 \\
 &  & Macro-F1 & 0.399 & 0.225 & 0.271 & 0.413 & 0.402 & 0.142 & 0.24 & 0.65 & 1.00 \\
 &  & RPS & 0.345 & 0.197 & 0.308 & 0.651 & 0.349 & 0.134 & 0.00 & 0.67 & 0.00 \\
 &  & QWK-floor & 0.443 & 0.232 & 0.249 & 0.302 & 0.412 & 0.146 & 0.45 & 0.72 & 1.00 \\
\addlinespace
SUPPORT2 & PO & QWK & 0.419 & 0.314 & 0.340 & 0.435 & 0.268 & 0.217 & 0.65 & 0.45 & 0.67 \\
 &  & MAE & 0.415 & 0.308 & 0.383 & 0.536 & 0.265 & 0.194 & 0.80 & 0.69 & 0.00 \\
 &  & AMAE & 0.419 & 0.318 & 0.312 & 0.405 & 0.259 & 0.228 & 0.58 & 0.34 & 1.00 \\
 &  & MMAE & 0.418 & 0.318 & 0.311 & 0.405 & 0.260 & 0.228 & 0.58 & 0.34 & 1.00 \\
 &  & Macro-F1 & 0.415 & 0.312 & 0.363 & 0.469 & 0.269 & 0.205 & 0.73 & 0.57 & 0.33 \\
 &  & RPS & 0.412 & 0.309 & 0.386 & 0.535 & 0.264 & 0.194 & 0.80 & 0.68 & 0.00 \\
 &  & QWK-floor & 0.419 & 0.314 & 0.340 & 0.435 & 0.268 & 0.217 & 0.65 & 0.45 & 0.67 \\
\addlinespace
 & RF & QWK & 0.431 & 0.300 & 0.367 & 0.529 & 0.298 & 0.193 & 0.76 & 0.72 & 0.17 \\
 &  & MAE & 0.431 & 0.300 & 0.367 & 0.529 & 0.298 & 0.193 & 0.76 & 0.72 & 0.17 \\
 &  & AMAE & 0.309 & 0.344 & 0.320 & 0.486 & 0.314 & 0.243 & 0.61 & 0.34 & 1.00 \\
 &  & MMAE & 0.415 & 0.302 & 0.333 & 0.390 & 0.349 & 0.212 & 0.69 & 0.59 & 1.00 \\
 &  & Macro-F1 & 0.401 & 0.308 & 0.339 & 0.395 & 0.342 & 0.212 & 0.70 & 0.58 & 1.00 \\
 &  & RPS & 0.435 & 0.299 & 0.370 & 0.548 & 0.288 & 0.190 & 0.77 & 0.74 & 0.00 \\
 &  & QWK-floor & 0.431 & 0.300 & 0.367 & 0.529 & 0.298 & 0.193 & 0.76 & 0.72 & 0.17 \\
\addlinespace
TAE (control) & PO & QWK & 0.291 & 0.343 & 0.340 & 0.490 & 0.447 & 0.222 & 0.68 & 0.33 & 0.50 \\
 &  & MAE & 0.291 & 0.343 & 0.340 & 0.490 & 0.447 & 0.222 & 0.68 & 0.33 & 0.50 \\
 &  & AMAE & 0.291 & 0.343 & 0.340 & 0.490 & 0.447 & 0.222 & 0.68 & 0.33 & 0.50 \\
 &  & MMAE & 0.300 & 0.353 & 0.350 & 0.495 & 0.421 & 0.221 & 0.76 & 0.35 & 0.50 \\
 &  & Macro-F1 & 0.291 & 0.343 & 0.340 & 0.490 & 0.447 & 0.222 & 0.68 & 0.33 & 0.67 \\
 &  & RPS & 0.247 & 0.384 & 0.381 & 0.529 & 0.377 & 0.224 & 0.69 & 0.35 & 0.67 \\
 &  & QWK-floor & 0.286 & 0.351 & 0.348 & 0.490 & 0.434 & 0.222 & 0.68 & 0.35 & 0.50 \\
\addlinespace
 & RF & QWK & 0.221 & 0.392 & 0.390 & 0.490 & 0.377 & 0.227 & 0.63 & 0.39 & 0.50 \\
 &  & MAE & 0.238 & 0.391 & 0.389 & 0.476 & 0.380 & 0.226 & 0.65 & 0.44 & 0.67 \\
 &  & AMAE & 0.238 & 0.391 & 0.389 & 0.476 & 0.380 & 0.226 & 0.65 & 0.44 & 0.67 \\
 &  & MMAE & 0.266 & 0.379 & 0.377 & 0.462 & 0.393 & 0.229 & 0.66 & 0.45 & 0.83 \\
 &  & Macro-F1 & 0.289 & 0.369 & 0.367 & 0.462 & 0.406 & 0.228 & 0.69 & 0.45 & 0.50 \\
 &  & RPS & 0.283 & 0.376 & 0.374 & 0.457 & 0.390 & 0.224 & 0.67 & 0.46 & 0.83 \\
 &  & QWK-floor & 0.216 & 0.401 & 0.399 & 0.476 & 0.367 & 0.227 & 0.62 & 0.44 & 0.83 \\
\addlinespace
Diabetes130 (control) & PO & QWK & 0.240 & 0.332 & 0.362 & 0.439 & 0.412 & 0.198 & 0.60 & 0.44 & 1.00 \\
 &  & MAE & 0.191 & 0.265 & 0.399 & 0.766 & 0.374 & 0.168 & 0.84 & 0.04 & 0.00 \\
 &  & AMAE & 0.240 & 0.332 & 0.362 & 0.439 & 0.412 & 0.198 & 0.60 & 0.44 & 1.00 \\
 &  & MMAE & 0.238 & 0.333 & 0.363 & 0.443 & 0.411 & 0.199 & 0.60 & 0.44 & 1.00 \\
 &  & Macro-F1 & 0.239 & 0.332 & 0.362 & 0.440 & 0.412 & 0.198 & 0.60 & 0.44 & 1.00 \\
 &  & RPS & 0.191 & 0.265 & 0.399 & 0.766 & 0.374 & 0.168 & 0.84 & 0.04 & 0.00 \\
 &  & QWK-floor & 0.240 & 0.332 & 0.362 & 0.439 & 0.412 & 0.198 & 0.60 & 0.44 & 1.00 \\
\addlinespace
 & RF & QWK & 0.228 & 0.347 & 0.375 & 0.438 & 0.394 & 0.206 & 0.65 & 0.46 & 1.00 \\
 &  & MAE & 0.161 & 0.265 & 0.405 & 0.794 & 0.358 & 0.170 & 0.83 & 0.00 & 0.00 \\
 &  & AMAE & 0.216 & 0.311 & 0.370 & 0.539 & 0.419 & 0.192 & 0.64 & 0.29 & 1.00 \\
 &  & MMAE & 0.228 & 0.347 & 0.375 & 0.438 & 0.394 & 0.206 & 0.65 & 0.46 & 1.00 \\
 &  & Macro-F1 & 0.211 & 0.292 & 0.366 & 0.590 & 0.421 & 0.184 & 0.64 & 0.20 & 1.00 \\
 &  & RPS & 0.161 & 0.265 & 0.405 & 0.794 & 0.358 & 0.170 & 0.83 & 0.00 & 0.00 \\
 &  & QWK-floor & 0.228 & 0.347 & 0.375 & 0.438 & 0.394 & 0.206 & 0.65 & 0.46 & 1.00 \\
\addlinespace
Student (control) & PO & QWK & 0.505 & 0.243 & 0.245 & 0.300 & 0.338 & 0.147 & 0.52 & 0.39 & 1.00 \\
 &  & MAE & 0.450 & 0.215 & 0.245 & 0.442 & 0.342 & 0.141 & 0.30 & 0.17 & 0.00 \\
 &  & AMAE & 0.496 & 0.240 & 0.247 & 0.326 & 0.336 & 0.146 & 0.48 & 0.35 & 0.83 \\
 &  & MMAE & 0.503 & 0.243 & 0.245 & 0.303 & 0.338 & 0.147 & 0.52 & 0.38 & 1.00 \\
 &  & Macro-F1 & 0.506 & 0.245 & 0.246 & 0.305 & 0.330 & 0.147 & 0.54 & 0.37 & 1.00 \\
 &  & RPS & 0.443 & 0.215 & 0.246 & 0.448 & 0.343 & 0.141 & 0.30 & 0.16 & 0.00 \\
 &  & QWK-floor & 0.505 & 0.243 & 0.245 & 0.300 & 0.338 & 0.147 & 0.52 & 0.39 & 1.00 \\
\addlinespace
 & RF & QWK & 0.462 & 0.259 & 0.255 & 0.323 & 0.310 & 0.146 & 0.67 & 0.27 & 1.00 \\
 &  & MAE & 0.308 & 0.230 & 0.273 & 0.535 & 0.268 & 0.146 & 0.20 & 0.05 & 0.00 \\
 &  & AMAE & 0.448 & 0.254 & 0.256 & 0.338 & 0.315 & 0.146 & 0.60 & 0.24 & 0.83 \\
 &  & MMAE & 0.468 & 0.258 & 0.254 & 0.309 & 0.309 & 0.147 & 0.66 & 0.28 & 1.00 \\
 &  & Macro-F1 & 0.448 & 0.254 & 0.256 & 0.338 & 0.315 & 0.146 & 0.60 & 0.24 & 0.83 \\
 &  & RPS & 0.424 & 0.234 & 0.255 & 0.402 & 0.325 & 0.143 & 0.39 & 0.15 & 0.33 \\
 &  & QWK-floor & 0.462 & 0.259 & 0.255 & 0.323 & 0.310 & 0.146 & 0.67 & 0.27 & 1.00 \\
\addlinespace
NPHA (control) & PO & QWK & 0.101 & 0.374 & 0.398 & 0.447 & 0.350 & 0.195 & 0.36 & 0.36 & 0.83 \\
 &  & MAE & 0.081 & 0.241 & 0.328 & 0.511 & 0.288 & 0.177 & 0.00 & 0.12 & 0.00 \\
 &  & AMAE & 0.089 & 0.240 & 0.327 & 0.513 & 0.293 & 0.177 & 0.00 & 0.13 & 0.00 \\
 &  & MMAE & 0.114 & 0.397 & 0.405 & 0.441 & 0.338 & 0.196 & 0.44 & 0.41 & 1.00 \\
 &  & Macro-F1 & 0.115 & 0.390 & 0.401 & 0.443 & 0.351 & 0.197 & 0.44 & 0.42 & 1.00 \\
 &  & RPS & 0.081 & 0.241 & 0.328 & 0.511 & 0.288 & 0.177 & 0.00 & 0.12 & 0.00 \\
 &  & QWK-floor & 0.110 & 0.395 & 0.405 & 0.443 & 0.343 & 0.197 & 0.44 & 0.41 & 1.00 \\
\addlinespace
 & RF & QWK & 0.071 & 0.387 & 0.412 & 0.461 & 0.347 & 0.190 & 0.34 & 0.39 & 1.00 \\
 &  & MAE & 0.034 & 0.241 & 0.333 & 0.515 & 0.265 & 0.177 & 0.00 & 0.06 & 0.00 \\
 &  & AMAE & 0.034 & 0.241 & 0.333 & 0.515 & 0.265 & 0.177 & 0.00 & 0.06 & 0.00 \\
 &  & MMAE & 0.069 & 0.375 & 0.405 & 0.479 & 0.346 & 0.186 & 0.32 & 0.34 & 0.83 \\
 &  & Macro-F1 & 0.091 & 0.376 & 0.403 & 0.468 & 0.355 & 0.190 & 0.34 & 0.40 & 1.00 \\
 &  & RPS & 0.034 & 0.241 & 0.333 & 0.515 & 0.265 & 0.177 & 0.00 & 0.06 & 0.00 \\
 &  & QWK-floor & 0.071 & 0.387 & 0.412 & 0.461 & 0.347 & 0.190 & 0.34 & 0.39 & 1.00 \\
\addlinespace
SkillCraft, 8 leagues (superseded) & PO & QWK & 0.736 & 0.125 & 0.108 & 0.138 & 0.368 & 0.085 & 0.58 & 0.75 & 1.00 \\
 &  & MAE & 0.724 & 0.108 & 0.130 & 0.183 & 0.347 & 0.075 & 0.23 & 0.46 & 0.00 \\
 &  & AMAE & 0.736 & 0.125 & 0.108 & 0.138 & 0.368 & 0.085 & 0.58 & 0.75 & 1.00 \\
 &  & MMAE & 0.736 & 0.125 & 0.108 & 0.138 & 0.368 & 0.085 & 0.58 & 0.75 & 1.00 \\
 &  & Macro-F1 & 0.733 & 0.122 & 0.112 & 0.137 & 0.372 & 0.083 & 0.52 & 0.71 & 0.83 \\
 &  & RPS & 0.725 & 0.108 & 0.130 & 0.182 & 0.348 & 0.075 & 0.23 & 0.46 & 0.00 \\
 &  & QWK-floor & 0.736 & 0.125 & 0.108 & 0.138 & 0.368 & 0.085 & 0.58 & 0.75 & 1.00 \\
\addlinespace
 & RF & QWK & 0.750 & 0.112 & 0.105 & 0.132 & 0.433 & 0.072 & 0.45 & 0.87 & 1.00 \\
 &  & MAE & 0.738 & 0.106 & 0.115 & 0.162 & 0.407 & 0.071 & 0.32 & 0.80 & 0.00 \\
 &  & AMAE & 0.737 & 0.122 & 0.105 & 0.146 & 0.403 & 0.077 & 0.54 & 0.87 & 1.00 \\
 &  & MMAE & 0.750 & 0.112 & 0.105 & 0.132 & 0.433 & 0.072 & 0.45 & 0.87 & 1.00 \\
 &  & Macro-F1 & 0.750 & 0.112 & 0.105 & 0.132 & 0.433 & 0.072 & 0.45 & 0.87 & 1.00 \\
 &  & RPS & 0.738 & 0.106 & 0.115 & 0.162 & 0.407 & 0.071 & 0.32 & 0.80 & 0.00 \\
 &  & QWK-floor & 0.750 & 0.112 & 0.105 & 0.132 & 0.433 & 0.072 & 0.45 & 0.87 & 1.00 \\
\addlinespace

\end{longtable}
}

{\footnotesize
\setlength{\tabcolsep}{3pt}
\begin{longtable}{@{}lllrlrrrrr@{}}
\caption{Differences relative to QWK selection. $\Delta R_{\mathrm{end}}$ with 95\% paired group-bootstrap interval (conditional on fitted models), the difference in each repeat, and differences in QWK and normalized MAE, AMAE and MMAE.}\label{tab:s-contrasts}\\
\toprule
Task & Fam. & Criterion & $\Delta R_{\mathrm{end}}$ & 95\% interval & Repeat 1 / 2 & $\Delta$QWK & $\Delta$MAE & $\Delta$AMAE & $\Delta$MMAE\\
\midrule\endfirsthead
\toprule
Task & Fam. & Criterion & $\Delta R_{\mathrm{end}}$ & 95\% interval & Repeat 1 / 2 & $\Delta$QWK & $\Delta$MAE & $\Delta$AMAE & $\Delta$MMAE\\
\midrule\endhead
\bottomrule\endfoot
Wine (red) & PO & AMAE & \ensuremath{+}0.531 & [\ensuremath{+}0.333, \ensuremath{+}0.730] & \ensuremath{+}0.506 / \ensuremath{+}0.556 & \ensuremath{-}0.034 & \ensuremath{+}0.089 & \ensuremath{-}0.059 & \ensuremath{-}0.217 \\
 &  & Macro-F1 & 0.000 & [0.000, 0.000] & 0.000 / 0.000 & 0.000 & 0.000 & 0.000 & 0.000 \\
 &  & QWK+floor & \ensuremath{+}0.531 & [\ensuremath{+}0.333, \ensuremath{+}0.750] & \ensuremath{+}0.506 / \ensuremath{+}0.556 & \ensuremath{-}0.031 & \ensuremath{+}0.088 & \ensuremath{-}0.059 & \ensuremath{-}0.218 \\
 &  & MMAE & \ensuremath{+}0.531 & [\ensuremath{+}0.333, \ensuremath{+}0.725] & \ensuremath{+}0.506 / \ensuremath{+}0.556 & \ensuremath{-}0.032 & \ensuremath{+}0.088 & \ensuremath{-}0.059 & \ensuremath{-}0.218 \\
 &  & MAE & 0.000 & [0.000, 0.000] & 0.000 / 0.000 & \ensuremath{-}0.003 & 0.000 & \ensuremath{+}0.002 & \ensuremath{+}0.010 \\
 &  & RPS & 0.000 & [0.000, 0.000] & 0.000 / 0.000 & \ensuremath{-}0.003 & 0.000 & 0.000 & 0.000 \\
\addlinespace
 & RF & AMAE & \ensuremath{+}0.108 & [\ensuremath{+}0.031, \ensuremath{+}0.203] & \ensuremath{+}0.083 / \ensuremath{+}0.133 & \ensuremath{-}0.056 & \ensuremath{+}0.037 & \ensuremath{-}0.005 & \ensuremath{-}0.050 \\
 &  & Macro-F1 & 0.000 & [0.000, 0.000] & 0.000 / 0.000 & 0.000 & 0.000 & 0.000 & 0.000 \\
 &  & QWK+floor & \ensuremath{+}0.108 & [\ensuremath{+}0.031, \ensuremath{+}0.202] & \ensuremath{+}0.083 / \ensuremath{+}0.133 & \ensuremath{-}0.056 & \ensuremath{+}0.037 & \ensuremath{-}0.005 & \ensuremath{-}0.050 \\
 &  & MMAE & \ensuremath{+}0.108 & [\ensuremath{+}0.033, \ensuremath{+}0.208] & \ensuremath{+}0.083 / \ensuremath{+}0.133 & \ensuremath{-}0.056 & \ensuremath{+}0.037 & \ensuremath{-}0.005 & \ensuremath{-}0.050 \\
 &  & MAE & \ensuremath{-}0.028 & [\ensuremath{-}0.074, 0.000] & \ensuremath{-}0.028 / \ensuremath{-}0.028 & \ensuremath{-}0.048 & \ensuremath{-}0.016 & \ensuremath{+}0.053 & \ensuremath{+}0.130 \\
 &  & RPS & \ensuremath{-}0.028 & [\ensuremath{-}0.068, 0.000] & \ensuremath{-}0.028 / \ensuremath{-}0.028 & \ensuremath{-}0.048 & \ensuremath{-}0.016 & \ensuremath{+}0.053 & \ensuremath{+}0.130 \\
\addlinespace
Wine (white) & PO & AMAE & \ensuremath{+}0.100 & [\ensuremath{+}0.012, \ensuremath{+}0.228] & \ensuremath{+}0.075 / \ensuremath{+}0.125 & \ensuremath{-}0.022 & \ensuremath{+}0.040 & \ensuremath{+}0.010 & \ensuremath{+}0.021 \\
 &  & Macro-F1 & 0.000 & [0.000, 0.000] & 0.000 / 0.000 & 0.000 & 0.000 & 0.000 & 0.000 \\
 &  & QWK+floor & \ensuremath{+}0.388 & [\ensuremath{+}0.216, \ensuremath{+}0.576] & \ensuremath{+}0.300 / \ensuremath{+}0.475 & \ensuremath{-}0.032 & \ensuremath{+}0.098 & \ensuremath{+}0.007 & \ensuremath{+}0.013 \\
 &  & MMAE & \ensuremath{+}0.037 & [0.000, \ensuremath{+}0.079] & \ensuremath{+}0.075 / 0.000 & \ensuremath{-}0.026 & \ensuremath{+}0.021 & \ensuremath{+}0.015 & 0.000 \\
 &  & MAE & 0.000 & [0.000, 0.000] & 0.000 / 0.000 & 0.000 & 0.000 & 0.000 & 0.000 \\
 &  & RPS & 0.000 & [0.000, 0.000] & 0.000 / 0.000 & 0.000 & 0.000 & 0.000 & 0.000 \\
\addlinespace
 & RF & AMAE & 0.000 & [0.000, 0.000] & 0.000 / 0.000 & 0.000 & 0.000 & 0.000 & 0.000 \\
 &  & Macro-F1 & 0.000 & [0.000, 0.000] & 0.000 / 0.000 & 0.000 & 0.000 & 0.000 & 0.000 \\
 &  & QWK+floor & 0.000 & [0.000, 0.000] & 0.000 / 0.000 & 0.000 & 0.000 & 0.000 & 0.000 \\
 &  & MMAE & \ensuremath{+}0.062 & [0.000, \ensuremath{+}0.181] & \ensuremath{+}0.025 / \ensuremath{+}0.100 & \ensuremath{-}0.100 & \ensuremath{+}0.032 & \ensuremath{+}0.015 & \ensuremath{-}0.004 \\
 &  & MAE & \ensuremath{-}0.037 & [\ensuremath{-}0.081, 0.000] & 0.000 / \ensuremath{-}0.075 & \ensuremath{-}0.090 & \ensuremath{-}0.011 & \ensuremath{+}0.034 & \ensuremath{+}0.050 \\
 &  & RPS & \ensuremath{-}0.037 & [\ensuremath{-}0.083, 0.000] & 0.000 / \ensuremath{-}0.075 & \ensuremath{-}0.066 & \ensuremath{-}0.013 & \ensuremath{+}0.032 & \ensuremath{+}0.046 \\
\addlinespace
SkillCraft & PO & AMAE & \ensuremath{-}0.001 & [\ensuremath{-}0.005, 0.000] & 0.000 / \ensuremath{-}0.003 & 0.000 & 0.000 & 0.000 & 0.000 \\
 &  & Macro-F1 & \ensuremath{-}0.004 & [\ensuremath{-}0.010, 0.000] & 0.000 / \ensuremath{-}0.009 & \ensuremath{-}0.001 & 0.000 & 0.000 & 0.000 \\
 &  & QWK+floor & 0.000 & [0.000, 0.000] & 0.000 / 0.000 & 0.000 & 0.000 & 0.000 & 0.000 \\
 &  & MMAE & 0.000 & [0.000, 0.000] & 0.000 / 0.000 & 0.000 & 0.000 & 0.000 & 0.000 \\
 &  & MAE & \ensuremath{-}0.497 & [\ensuremath{-}0.581, \ensuremath{-}0.412] & \ensuremath{-}0.493 / \ensuremath{-}0.502 & \ensuremath{-}0.020 & \ensuremath{-}0.015 & \ensuremath{+}0.022 & \ensuremath{+}0.057 \\
 &  & RPS & \ensuremath{-}0.497 & [\ensuremath{-}0.581, \ensuremath{-}0.412] & \ensuremath{-}0.493 / \ensuremath{-}0.502 & \ensuremath{-}0.020 & \ensuremath{-}0.015 & \ensuremath{+}0.022 & \ensuremath{+}0.057 \\
\addlinespace
 & RF & AMAE & \ensuremath{+}0.187 & [\ensuremath{+}0.133, \ensuremath{+}0.245] & \ensuremath{+}0.113 / \ensuremath{+}0.262 & \ensuremath{-}0.020 & \ensuremath{+}0.014 & \ensuremath{+}0.001 & \ensuremath{+}0.023 \\
 &  & Macro-F1 & \ensuremath{+}0.055 & [\ensuremath{+}0.024, \ensuremath{+}0.091] & 0.000 / \ensuremath{+}0.110 & \ensuremath{-}0.004 & \ensuremath{+}0.003 & \ensuremath{-}0.001 & \ensuremath{+}0.004 \\
 &  & QWK+floor & \ensuremath{+}0.249 & [\ensuremath{+}0.185, \ensuremath{+}0.323] & \ensuremath{+}0.237 / \ensuremath{+}0.262 & \ensuremath{-}0.023 & \ensuremath{+}0.018 & 0.000 & \ensuremath{+}0.028 \\
 &  & MMAE & 0.000 & [0.000, 0.000] & 0.000 / 0.000 & 0.000 & 0.000 & 0.000 & 0.000 \\
 &  & MAE & \ensuremath{-}0.163 & [\ensuremath{-}0.221, \ensuremath{-}0.113] & \ensuremath{-}0.155 / \ensuremath{-}0.170 & \ensuremath{-}0.025 & \ensuremath{-}0.006 & \ensuremath{+}0.013 & \ensuremath{+}0.052 \\
 &  & RPS & \ensuremath{-}0.146 & [\ensuremath{-}0.205, \ensuremath{-}0.094] & \ensuremath{-}0.137 / \ensuremath{-}0.155 & \ensuremath{-}0.019 & \ensuremath{-}0.007 & \ensuremath{+}0.011 & \ensuremath{+}0.039 \\
\addlinespace
Restaurant & PO & AMAE & \ensuremath{-}0.008 & [\ensuremath{-}0.026, \ensuremath{+}0.008] & \ensuremath{-}0.016 / 0.000 & \ensuremath{+}0.008 & \ensuremath{-}0.017 & \ensuremath{-}0.009 & \ensuremath{+}0.035 \\
 &  & Macro-F1 & \ensuremath{+}0.026 & [\ensuremath{-}0.006, \ensuremath{+}0.063] & \ensuremath{+}0.053 / 0.000 & \ensuremath{-}0.003 & \ensuremath{+}0.019 & \ensuremath{+}0.006 & \ensuremath{-}0.069 \\
 &  & QWK+floor & \ensuremath{+}0.047 & [\ensuremath{+}0.005, \ensuremath{+}0.092] & \ensuremath{+}0.064 / \ensuremath{+}0.031 & \ensuremath{-}0.029 & \ensuremath{+}0.065 & \ensuremath{+}0.029 & \ensuremath{-}0.157 \\
 &  & MMAE & \ensuremath{+}0.049 & [\ensuremath{+}0.006, \ensuremath{+}0.093] & \ensuremath{+}0.067 / \ensuremath{+}0.032 & \ensuremath{-}0.023 & \ensuremath{+}0.062 & \ensuremath{+}0.027 & \ensuremath{-}0.162 \\
 &  & MAE & \ensuremath{-}0.022 & [\ensuremath{-}0.056, \ensuremath{+}0.008] & \ensuremath{-}0.016 / \ensuremath{-}0.028 & \ensuremath{+}0.005 & \ensuremath{-}0.037 & \ensuremath{-}0.017 & \ensuremath{+}0.094 \\
 &  & RPS & \ensuremath{-}0.022 & [\ensuremath{-}0.057, \ensuremath{+}0.008] & \ensuremath{-}0.016 / \ensuremath{-}0.028 & \ensuremath{+}0.005 & \ensuremath{-}0.037 & \ensuremath{-}0.017 & \ensuremath{+}0.094 \\
\addlinespace
 & RF & AMAE & \ensuremath{-}0.006 & [\ensuremath{-}0.017, \ensuremath{+}0.006] & \ensuremath{-}0.012 / 0.000 & \ensuremath{-}0.010 & \ensuremath{+}0.005 & \ensuremath{+}0.002 & \ensuremath{-}0.016 \\
 &  & Macro-F1 & \ensuremath{+}0.003 & [\ensuremath{-}0.014, \ensuremath{+}0.023] & \ensuremath{-}0.012 / \ensuremath{+}0.018 & \ensuremath{-}0.002 & \ensuremath{+}0.010 & 0.000 & \ensuremath{-}0.048 \\
 &  & QWK+floor & \ensuremath{+}0.001 & [\ensuremath{-}0.041, \ensuremath{+}0.040] & \ensuremath{-}0.004 / \ensuremath{+}0.006 & \ensuremath{-}0.033 & \ensuremath{+}0.046 & \ensuremath{+}0.012 & \ensuremath{-}0.161 \\
 &  & MMAE & \ensuremath{+}0.004 & [\ensuremath{-}0.038, \ensuremath{+}0.047] & \ensuremath{+}0.002 / \ensuremath{+}0.006 & \ensuremath{-}0.020 & \ensuremath{+}0.045 & \ensuremath{+}0.007 & \ensuremath{-}0.175 \\
 &  & MAE & \ensuremath{+}0.006 & [\ensuremath{-}0.018, \ensuremath{+}0.026] & \ensuremath{+}0.008 / \ensuremath{+}0.003 & \ensuremath{+}0.007 & \ensuremath{-}0.023 & \ensuremath{-}0.005 & \ensuremath{+}0.089 \\
 &  & RPS & \ensuremath{+}0.015 & [\ensuremath{-}0.012, \ensuremath{+}0.039] & \ensuremath{+}0.008 / \ensuremath{+}0.023 & \ensuremath{+}0.019 & \ensuremath{-}0.032 & \ensuremath{-}0.008 & \ensuremath{+}0.114 \\
\addlinespace
Sushi & PO & AMAE & 0.000 & [0.000, 0.000] & 0.000 / 0.000 & 0.000 & 0.000 & 0.000 & 0.000 \\
 &  & Macro-F1 & \ensuremath{+}0.003 & [0.000, \ensuremath{+}0.007] & \ensuremath{+}0.006 / \ensuremath{+}0.001 & 0.000 & \ensuremath{+}0.002 & \ensuremath{+}0.001 & \ensuremath{+}0.002 \\
 &  & QWK+floor & 0.000 & [0.000, 0.000] & 0.000 / 0.000 & 0.000 & 0.000 & 0.000 & 0.000 \\
 &  & MMAE & 0.000 & [0.000, \ensuremath{+}0.001] & \ensuremath{+}0.001 / 0.000 & 0.000 & \ensuremath{-}0.001 & 0.000 & \ensuremath{-}0.001 \\
 &  & MAE & \ensuremath{-}0.047 & [\ensuremath{-}0.077, \ensuremath{-}0.019] & \ensuremath{-}0.022 / \ensuremath{-}0.072 & \ensuremath{-}0.060 & \ensuremath{-}0.094 & \ensuremath{+}0.036 & \ensuremath{+}0.404 \\
 &  & RPS & \ensuremath{-}0.035 & [\ensuremath{-}0.066, \ensuremath{-}0.005] & \ensuremath{-}0.015 / \ensuremath{-}0.056 & \ensuremath{-}0.071 & \ensuremath{-}0.090 & \ensuremath{+}0.045 & \ensuremath{+}0.421 \\
\addlinespace
 & RF & AMAE & \ensuremath{-}0.005 & [\ensuremath{-}0.010, 0.000] & 0.000 / \ensuremath{-}0.010 & \ensuremath{+}0.007 & \ensuremath{-}0.004 & \ensuremath{-}0.006 & 0.000 \\
 &  & Macro-F1 & \ensuremath{-}0.005 & [\ensuremath{-}0.010, 0.000] & 0.000 / \ensuremath{-}0.010 & \ensuremath{+}0.007 & \ensuremath{-}0.004 & \ensuremath{-}0.006 & 0.000 \\
 &  & QWK+floor & 0.000 & [0.000, 0.000] & 0.000 / 0.000 & 0.000 & 0.000 & 0.000 & 0.000 \\
 &  & MMAE & \ensuremath{-}0.002 & [\ensuremath{-}0.007, \ensuremath{+}0.003] & \ensuremath{-}0.003 / 0.000 & \ensuremath{-}0.001 & 0.000 & \ensuremath{+}0.002 & \ensuremath{-}0.004 \\
 &  & MAE & \ensuremath{+}0.061 & [\ensuremath{+}0.034, \ensuremath{+}0.087] & \ensuremath{+}0.081 / \ensuremath{+}0.040 & \ensuremath{-}0.071 & \ensuremath{-}0.066 & \ensuremath{+}0.071 & \ensuremath{+}0.399 \\
 &  & RPS & \ensuremath{+}0.142 & [\ensuremath{+}0.114, \ensuremath{+}0.167] & \ensuremath{+}0.171 / \ensuremath{+}0.113 & \ensuremath{-}0.128 & \ensuremath{-}0.047 & \ensuremath{+}0.126 & \ensuremath{+}0.522 \\
\addlinespace
MovieLens-100k & PO & AMAE & \ensuremath{-}0.403 & [\ensuremath{-}0.427, \ensuremath{-}0.382] & \ensuremath{-}0.403 / \ensuremath{-}0.404 & \ensuremath{-}0.060 & \ensuremath{-}0.149 & \ensuremath{-}0.031 & \ensuremath{+}0.281 \\
 &  & Macro-F1 & 0.000 & [0.000, 0.000] & 0.000 / 0.000 & 0.000 & 0.000 & 0.000 & 0.000 \\
 &  & QWK+floor & 0.000 & [0.000, 0.000] & 0.000 / 0.000 & 0.000 & 0.000 & 0.000 & 0.000 \\
 &  & MMAE & 0.000 & [0.000, 0.000] & 0.000 / 0.000 & 0.000 & 0.000 & 0.000 & 0.000 \\
 &  & MAE & \ensuremath{-}0.403 & [\ensuremath{-}0.428, \ensuremath{-}0.383] & \ensuremath{-}0.403 / \ensuremath{-}0.404 & \ensuremath{-}0.060 & \ensuremath{-}0.149 & \ensuremath{-}0.031 & \ensuremath{+}0.281 \\
 &  & RPS & \ensuremath{-}0.403 & [\ensuremath{-}0.428, \ensuremath{-}0.383] & \ensuremath{-}0.403 / \ensuremath{-}0.404 & \ensuremath{-}0.060 & \ensuremath{-}0.149 & \ensuremath{-}0.031 & \ensuremath{+}0.281 \\
\addlinespace
 & RF & AMAE & \ensuremath{-}0.091 & [\ensuremath{-}0.101, \ensuremath{-}0.082] & \ensuremath{-}0.086 / \ensuremath{-}0.095 & \ensuremath{-}0.009 & \ensuremath{-}0.046 & \ensuremath{-}0.004 & \ensuremath{+}0.128 \\
 &  & Macro-F1 & \ensuremath{-}0.040 & [\ensuremath{-}0.052, \ensuremath{-}0.027] & \ensuremath{-}0.040 / \ensuremath{-}0.041 & \ensuremath{-}0.003 & \ensuremath{-}0.023 & \ensuremath{-}0.003 & \ensuremath{+}0.058 \\
 &  & QWK+floor & \ensuremath{+}0.050 & [\ensuremath{+}0.034, \ensuremath{+}0.069] & \ensuremath{+}0.054 / \ensuremath{+}0.045 & \ensuremath{-}0.017 & \ensuremath{+}0.026 & \ensuremath{+}0.018 & \ensuremath{-}0.009 \\
 &  & MMAE & \ensuremath{+}0.077 & [\ensuremath{+}0.054, \ensuremath{+}0.105] & \ensuremath{+}0.080 / \ensuremath{+}0.074 & \ensuremath{-}0.019 & \ensuremath{+}0.043 & \ensuremath{+}0.026 & \ensuremath{-}0.030 \\
 &  & MAE & \ensuremath{-}0.310 & [\ensuremath{-}0.338, \ensuremath{-}0.286] & \ensuremath{-}0.308 / \ensuremath{-}0.312 & \ensuremath{-}0.195 & \ensuremath{-}0.088 & \ensuremath{+}0.022 & \ensuremath{+}0.346 \\
 &  & RPS & \ensuremath{-}0.310 & [\ensuremath{-}0.338, \ensuremath{-}0.286] & \ensuremath{-}0.308 / \ensuremath{-}0.312 & \ensuremath{-}0.195 & \ensuremath{-}0.088 & \ensuremath{+}0.022 & \ensuremath{+}0.346 \\
\addlinespace
CTG & PO & AMAE & \ensuremath{+}0.034 & [\ensuremath{-}0.003, \ensuremath{+}0.070] & \ensuremath{+}0.052 / \ensuremath{+}0.016 & \ensuremath{-}0.026 & \ensuremath{+}0.024 & \ensuremath{-}0.052 & \ensuremath{-}0.108 \\
 &  & Macro-F1 & \ensuremath{+}0.013 & [\ensuremath{-}0.016, \ensuremath{+}0.036] & \ensuremath{+}0.042 / \ensuremath{-}0.016 & \ensuremath{-}0.038 & \ensuremath{+}0.024 & \ensuremath{-}0.036 & \ensuremath{-}0.069 \\
 &  & QWK+floor & 0.000 & [0.000, 0.000] & 0.000 / 0.000 & 0.000 & 0.000 & 0.000 & 0.000 \\
 &  & MMAE & \ensuremath{+}0.034 & [\ensuremath{-}0.005, \ensuremath{+}0.070] & \ensuremath{+}0.052 / \ensuremath{+}0.016 & \ensuremath{-}0.026 & \ensuremath{+}0.024 & \ensuremath{-}0.052 & \ensuremath{-}0.108 \\
 &  & MAE & \ensuremath{-}0.012 & [\ensuremath{-}0.041, \ensuremath{+}0.009] & \ensuremath{-}0.004 / \ensuremath{-}0.020 & 0.000 & \ensuremath{-}0.005 & \ensuremath{+}0.019 & \ensuremath{+}0.034 \\
 &  & RPS & \ensuremath{-}0.002 & [\ensuremath{-}0.035, \ensuremath{+}0.029] & \ensuremath{+}0.017 / \ensuremath{-}0.020 & \ensuremath{+}0.005 & \ensuremath{-}0.006 & \ensuremath{+}0.015 & \ensuremath{+}0.031 \\
\addlinespace
 & RF & AMAE & \ensuremath{+}0.008 & [\ensuremath{-}0.014, \ensuremath{+}0.035] & \ensuremath{+}0.014 / \ensuremath{+}0.002 & \ensuremath{-}0.041 & \ensuremath{+}0.014 & \ensuremath{-}0.026 & \ensuremath{-}0.054 \\
 &  & Macro-F1 & \ensuremath{+}0.012 & [\ensuremath{-}0.003, \ensuremath{+}0.028] & \ensuremath{+}0.006 / \ensuremath{+}0.018 & \ensuremath{-}0.021 & \ensuremath{+}0.006 & \ensuremath{-}0.015 & \ensuremath{-}0.032 \\
 &  & QWK+floor & 0.000 & [0.000, 0.000] & 0.000 / 0.000 & 0.000 & 0.000 & 0.000 & 0.000 \\
 &  & MMAE & \ensuremath{+}0.004 & [\ensuremath{-}0.019, \ensuremath{+}0.030] & \ensuremath{+}0.005 / \ensuremath{+}0.002 & \ensuremath{-}0.046 & \ensuremath{+}0.016 & \ensuremath{-}0.025 & \ensuremath{-}0.056 \\
 &  & MAE & 0.000 & [0.000, 0.000] & 0.000 / 0.000 & 0.000 & 0.000 & 0.000 & 0.000 \\
 &  & RPS & 0.000 & [0.000, 0.000] & 0.000 / 0.000 & 0.000 & 0.000 & 0.000 & 0.000 \\
\addlinespace
HCV & PO & AMAE & \ensuremath{-}0.142 & [\ensuremath{-}0.197, \ensuremath{-}0.080] & \ensuremath{-}0.151 / \ensuremath{-}0.134 & \ensuremath{-}0.106 & \ensuremath{+}0.084 & \ensuremath{-}0.113 & \ensuremath{-}0.211 \\
 &  & Macro-F1 & \ensuremath{-}0.043 & [\ensuremath{-}0.070, \ensuremath{-}0.013] & \ensuremath{-}0.031 / \ensuremath{-}0.056 & \ensuremath{-}0.061 & \ensuremath{+}0.036 & \ensuremath{-}0.044 & \ensuremath{-}0.103 \\
 &  & QWK+floor & 0.000 & [0.000, 0.000] & 0.000 / 0.000 & 0.000 & 0.000 & 0.000 & 0.000 \\
 &  & MMAE & \ensuremath{-}0.142 & [\ensuremath{-}0.194, \ensuremath{-}0.080] & \ensuremath{-}0.151 / \ensuremath{-}0.134 & \ensuremath{-}0.106 & \ensuremath{+}0.084 & \ensuremath{-}0.113 & \ensuremath{-}0.211 \\
 &  & MAE & 0.000 & [0.000, 0.000] & 0.000 / 0.000 & \ensuremath{-}0.014 & \ensuremath{+}0.001 & \ensuremath{+}0.006 & 0.000 \\
 &  & RPS & \ensuremath{-}0.008 & [\ensuremath{-}0.039, \ensuremath{+}0.020] & 0.000 / \ensuremath{-}0.016 & \ensuremath{-}0.050 & \ensuremath{+}0.005 & \ensuremath{+}0.034 & \ensuremath{+}0.087 \\
\addlinespace
 & RF & AMAE & \ensuremath{-}0.051 & [\ensuremath{-}0.087, \ensuremath{-}0.024] & \ensuremath{-}0.026 / \ensuremath{-}0.076 & \ensuremath{-}0.055 & \ensuremath{+}0.015 & \ensuremath{-}0.008 & \ensuremath{-}0.038 \\
 &  & Macro-F1 & \ensuremath{-}0.017 & [\ensuremath{-}0.037, \ensuremath{-}0.006] & \ensuremath{-}0.026 / \ensuremath{-}0.008 & \ensuremath{-}0.019 & \ensuremath{+}0.005 & \ensuremath{-}0.011 & \ensuremath{-}0.032 \\
 &  & QWK+floor & 0.000 & [0.000, 0.000] & 0.000 / 0.000 & 0.000 & 0.000 & 0.000 & 0.000 \\
 &  & MMAE & \ensuremath{-}0.021 & [\ensuremath{-}0.043, \ensuremath{-}0.009] & \ensuremath{-}0.029 / \ensuremath{-}0.014 & \ensuremath{-}0.026 & \ensuremath{+}0.008 & \ensuremath{-}0.016 & \ensuremath{-}0.054 \\
 &  & MAE & \ensuremath{-}0.007 & [\ensuremath{-}0.025, \ensuremath{+}0.003] & \ensuremath{-}0.017 / \ensuremath{+}0.003 & \ensuremath{-}0.024 & \ensuremath{+}0.003 & \ensuremath{+}0.025 & \ensuremath{+}0.040 \\
 &  & RPS & \ensuremath{-}0.041 & [\ensuremath{-}0.088, \ensuremath{-}0.003] & \ensuremath{-}0.025 / \ensuremath{-}0.057 & \ensuremath{-}0.051 & \ensuremath{+}0.004 & \ensuremath{+}0.073 & \ensuremath{+}0.079 \\
\addlinespace
Cirrhosis & PO & AMAE & \ensuremath{+}0.039 & [\ensuremath{-}0.001, \ensuremath{+}0.082] & 0.000 / \ensuremath{+}0.078 & \ensuremath{-}0.001 & \ensuremath{+}0.017 & \ensuremath{+}0.007 & \ensuremath{-}0.024 \\
 &  & Macro-F1 & \ensuremath{+}0.008 & [\ensuremath{-}0.057, \ensuremath{+}0.071] & \ensuremath{-}0.061 / \ensuremath{+}0.078 & \ensuremath{+}0.002 & \ensuremath{+}0.008 & \ensuremath{+}0.013 & \ensuremath{+}0.024 \\
 &  & QWK+floor & \ensuremath{+}0.029 & [\ensuremath{-}0.007, \ensuremath{+}0.074] & 0.000 / \ensuremath{+}0.058 & \ensuremath{+}0.001 & \ensuremath{+}0.016 & \ensuremath{+}0.005 & \ensuremath{-}0.026 \\
 &  & MMAE & \ensuremath{+}0.015 & [\ensuremath{-}0.029, \ensuremath{+}0.062] & \ensuremath{-}0.027 / \ensuremath{+}0.058 & \ensuremath{-}0.002 & \ensuremath{+}0.018 & \ensuremath{+}0.007 & \ensuremath{-}0.016 \\
 &  & MAE & \ensuremath{-}0.156 & [\ensuremath{-}0.248, \ensuremath{-}0.078] & \ensuremath{-}0.180 / \ensuremath{-}0.132 & \ensuremath{-}0.005 & \ensuremath{-}0.064 & \ensuremath{+}0.007 & \ensuremath{+}0.230 \\
 &  & RPS & \ensuremath{-}0.158 & [\ensuremath{-}0.249, \ensuremath{-}0.073] & \ensuremath{-}0.184 / \ensuremath{-}0.132 & \ensuremath{-}0.002 & \ensuremath{-}0.066 & \ensuremath{+}0.006 & \ensuremath{+}0.230 \\
\addlinespace
 & RF & AMAE & \ensuremath{+}0.029 & [\ensuremath{-}0.002, \ensuremath{+}0.066] & 0.000 / \ensuremath{+}0.058 & \ensuremath{+}0.005 & \ensuremath{-}0.002 & \ensuremath{-}0.006 & \ensuremath{-}0.024 \\
 &  & Macro-F1 & \ensuremath{-}0.068 & [\ensuremath{-}0.143, \ensuremath{-}0.014] & \ensuremath{-}0.068 / \ensuremath{-}0.068 & \ensuremath{-}0.027 & \ensuremath{-}0.006 & \ensuremath{+}0.007 & \ensuremath{+}0.056 \\
 &  & QWK+floor & \ensuremath{+}0.075 & [\ensuremath{+}0.014, \ensuremath{+}0.148] & \ensuremath{+}0.089 / \ensuremath{+}0.062 & \ensuremath{+}0.018 & \ensuremath{+}0.001 & \ensuremath{-}0.015 & \ensuremath{-}0.055 \\
 &  & MMAE & \ensuremath{+}0.073 & [\ensuremath{+}0.013, \ensuremath{+}0.154] & \ensuremath{+}0.089 / \ensuremath{+}0.058 & \ensuremath{+}0.015 & \ensuremath{-}0.001 & \ensuremath{-}0.013 & \ensuremath{-}0.039 \\
 &  & MAE & \ensuremath{-}0.155 & [\ensuremath{-}0.263, \ensuremath{-}0.064] & \ensuremath{-}0.136 / \ensuremath{-}0.174 & \ensuremath{-}0.075 & \ensuremath{-}0.034 & \ensuremath{+}0.037 & \ensuremath{+}0.254 \\
 &  & RPS & \ensuremath{-}0.179 & [\ensuremath{-}0.289, \ensuremath{-}0.072] & \ensuremath{-}0.181 / \ensuremath{-}0.177 & \ensuremath{-}0.081 & \ensuremath{-}0.035 & \ensuremath{+}0.044 & \ensuremath{+}0.294 \\
\addlinespace
SUPPORT2 & PO & AMAE & \ensuremath{-}0.093 & [\ensuremath{-}0.099, \ensuremath{-}0.086] & \ensuremath{-}0.188 / \ensuremath{+}0.003 & \ensuremath{-}0.001 & \ensuremath{+}0.004 & \ensuremath{-}0.028 & \ensuremath{-}0.030 \\
 &  & Macro-F1 & \ensuremath{+}0.095 & [\ensuremath{+}0.088, \ensuremath{+}0.102] & \ensuremath{+}0.093 / \ensuremath{+}0.097 & \ensuremath{-}0.005 & \ensuremath{-}0.002 & \ensuremath{+}0.023 & \ensuremath{+}0.034 \\
 &  & QWK+floor & 0.000 & [0.000, 0.000] & 0.000 / 0.000 & 0.000 & 0.000 & 0.000 & 0.000 \\
 &  & MMAE & \ensuremath{-}0.093 & [\ensuremath{-}0.099, \ensuremath{-}0.086] & \ensuremath{-}0.188 / \ensuremath{+}0.003 & \ensuremath{-}0.001 & \ensuremath{+}0.004 & \ensuremath{-}0.028 & \ensuremath{-}0.030 \\
 &  & MAE & \ensuremath{+}0.188 & [\ensuremath{+}0.179, \ensuremath{+}0.199] & \ensuremath{+}0.093 / \ensuremath{+}0.283 & \ensuremath{-}0.004 & \ensuremath{-}0.005 & \ensuremath{+}0.044 & \ensuremath{+}0.101 \\
 &  & RPS & \ensuremath{+}0.187 & [\ensuremath{+}0.177, \ensuremath{+}0.196] & \ensuremath{+}0.091 / \ensuremath{+}0.283 & \ensuremath{-}0.007 & \ensuremath{-}0.004 & \ensuremath{+}0.046 & \ensuremath{+}0.100 \\
\addlinespace
 & RF & AMAE & \ensuremath{-}0.264 & [\ensuremath{-}0.278, \ensuremath{-}0.251] & \ensuremath{-}0.248 / \ensuremath{-}0.280 & \ensuremath{-}0.123 & \ensuremath{+}0.044 & \ensuremath{-}0.047 & \ensuremath{-}0.042 \\
 &  & Macro-F1 & \ensuremath{-}0.097 & [\ensuremath{-}0.106, \ensuremath{-}0.088] & \ensuremath{-}0.054 / \ensuremath{-}0.140 & \ensuremath{-}0.030 & \ensuremath{+}0.008 & \ensuremath{-}0.028 & \ensuremath{-}0.134 \\
 &  & QWK+floor & 0.000 & [0.000, 0.000] & 0.000 / 0.000 & 0.000 & 0.000 & 0.000 & 0.000 \\
 &  & MMAE & \ensuremath{-}0.094 & [\ensuremath{-}0.103, \ensuremath{-}0.085] & \ensuremath{-}0.115 / \ensuremath{-}0.073 & \ensuremath{-}0.016 & \ensuremath{+}0.002 & \ensuremath{-}0.034 & \ensuremath{-}0.139 \\
 &  & MAE & 0.000 & [0.000, 0.000] & 0.000 / 0.000 & 0.000 & 0.000 & 0.000 & 0.000 \\
 &  & RPS & \ensuremath{+}0.013 & [\ensuremath{+}0.010, \ensuremath{+}0.017] & \ensuremath{+}0.027 / 0.000 & \ensuremath{+}0.004 & \ensuremath{-}0.001 & \ensuremath{+}0.003 & \ensuremath{+}0.019 \\
\addlinespace
TAE (control) & PO & AMAE & 0.000 & [0.000, 0.000] & 0.000 / 0.000 & 0.000 & 0.000 & 0.000 & 0.000 \\
 &  & Macro-F1 & 0.000 & [0.000, 0.000] & 0.000 / 0.000 & 0.000 & 0.000 & 0.000 & 0.000 \\
 &  & QWK+floor & \ensuremath{+}0.010 & [0.000, \ensuremath{+}0.034] & 0.000 / \ensuremath{+}0.019 & \ensuremath{-}0.005 & \ensuremath{+}0.008 & \ensuremath{+}0.008 & 0.000 \\
 &  & MMAE & \ensuremath{+}0.045 & [0.000, \ensuremath{+}0.090] & 0.000 / \ensuremath{+}0.091 & \ensuremath{+}0.010 & \ensuremath{+}0.010 & \ensuremath{+}0.010 & \ensuremath{+}0.005 \\
 &  & MAE & 0.000 & [0.000, 0.000] & 0.000 / 0.000 & 0.000 & 0.000 & 0.000 & 0.000 \\
 &  & RPS & \ensuremath{+}0.015 & [\ensuremath{-}0.044, \ensuremath{+}0.068] & \ensuremath{-}0.061 / \ensuremath{+}0.091 & \ensuremath{-}0.044 & \ensuremath{+}0.041 & \ensuremath{+}0.041 & \ensuremath{+}0.038 \\
\addlinespace
 & RF & AMAE & \ensuremath{+}0.034 & [0.000, \ensuremath{+}0.080] & 0.000 / \ensuremath{+}0.068 & \ensuremath{+}0.017 & \ensuremath{-}0.002 & \ensuremath{-}0.002 & \ensuremath{-}0.014 \\
 &  & Macro-F1 & \ensuremath{+}0.059 & [\ensuremath{+}0.007, \ensuremath{+}0.114] & \ensuremath{-}0.010 / \ensuremath{+}0.129 & \ensuremath{+}0.068 & \ensuremath{-}0.023 & \ensuremath{-}0.023 & \ensuremath{-}0.029 \\
 &  & QWK+floor & \ensuremath{+}0.019 & [\ensuremath{-}0.019, \ensuremath{+}0.070] & \ensuremath{-}0.031 / \ensuremath{+}0.068 & \ensuremath{-}0.005 & \ensuremath{+}0.008 & \ensuremath{+}0.009 & \ensuremath{-}0.014 \\
 &  & MMAE & \ensuremath{+}0.044 & [\ensuremath{+}0.001, \ensuremath{+}0.094] & \ensuremath{-}0.040 / \ensuremath{+}0.129 & \ensuremath{+}0.045 & \ensuremath{-}0.013 & \ensuremath{-}0.013 & \ensuremath{-}0.029 \\
 &  & MAE & \ensuremath{+}0.034 & [0.000, \ensuremath{+}0.082] & 0.000 / \ensuremath{+}0.068 & \ensuremath{+}0.017 & \ensuremath{-}0.002 & \ensuremath{-}0.002 & \ensuremath{-}0.014 \\
 &  & RPS & \ensuremath{+}0.054 & [\ensuremath{+}0.005, \ensuremath{+}0.116] & \ensuremath{-}0.010 / \ensuremath{+}0.118 & \ensuremath{+}0.062 & \ensuremath{-}0.017 & \ensuremath{-}0.016 & \ensuremath{-}0.034 \\
\addlinespace
Diabetes130 (control) & PO & AMAE & 0.000 & [0.000, 0.000] & 0.000 / 0.000 & 0.000 & 0.000 & 0.000 & 0.000 \\
 &  & Macro-F1 & \ensuremath{-}0.002 & [\ensuremath{-}0.005, \ensuremath{+}0.001] & 0.000 / \ensuremath{-}0.004 & \ensuremath{-}0.001 & 0.000 & 0.000 & \ensuremath{+}0.002 \\
 &  & QWK+floor & 0.000 & [0.000, 0.000] & 0.000 / 0.000 & 0.000 & 0.000 & 0.000 & 0.000 \\
 &  & MMAE & \ensuremath{-}0.003 & [\ensuremath{-}0.006, \ensuremath{-}0.001] & \ensuremath{-}0.006 / 0.000 & \ensuremath{-}0.002 & \ensuremath{+}0.001 & \ensuremath{+}0.001 & \ensuremath{+}0.004 \\
 &  & MAE & \ensuremath{-}0.078 & [\ensuremath{-}0.098, \ensuremath{-}0.057] & \ensuremath{-}0.077 / \ensuremath{-}0.078 & \ensuremath{-}0.049 & \ensuremath{-}0.067 & \ensuremath{+}0.037 & \ensuremath{+}0.327 \\
 &  & RPS & \ensuremath{-}0.078 & [\ensuremath{-}0.100, \ensuremath{-}0.057] & \ensuremath{-}0.077 / \ensuremath{-}0.078 & \ensuremath{-}0.049 & \ensuremath{-}0.067 & \ensuremath{+}0.037 & \ensuremath{+}0.327 \\
\addlinespace
 & RF & AMAE & \ensuremath{-}0.094 & [\ensuremath{-}0.109, \ensuremath{-}0.080] & \ensuremath{-}0.079 / \ensuremath{-}0.109 & \ensuremath{-}0.013 & \ensuremath{-}0.035 & \ensuremath{-}0.005 & \ensuremath{+}0.101 \\
 &  & Macro-F1 & \ensuremath{-}0.141 & [\ensuremath{-}0.159, \ensuremath{-}0.123] & \ensuremath{-}0.115 / \ensuremath{-}0.166 & \ensuremath{-}0.017 & \ensuremath{-}0.055 & \ensuremath{-}0.009 & \ensuremath{+}0.152 \\
 &  & QWK+floor & 0.000 & [0.000, 0.000] & 0.000 / 0.000 & 0.000 & 0.000 & 0.000 & 0.000 \\
 &  & MMAE & 0.000 & [0.000, 0.000] & 0.000 / 0.000 & 0.000 & 0.000 & 0.000 & 0.000 \\
 &  & MAE & \ensuremath{-}0.140 & [\ensuremath{-}0.161, \ensuremath{-}0.119] & \ensuremath{-}0.127 / \ensuremath{-}0.152 & \ensuremath{-}0.067 & \ensuremath{-}0.081 & \ensuremath{+}0.029 & \ensuremath{+}0.355 \\
 &  & RPS & \ensuremath{-}0.140 & [\ensuremath{-}0.161, \ensuremath{-}0.119] & \ensuremath{-}0.127 / \ensuremath{-}0.152 & \ensuremath{-}0.067 & \ensuremath{-}0.081 & \ensuremath{+}0.029 & \ensuremath{+}0.355 \\
\addlinespace
Student (control) & PO & AMAE & \ensuremath{-}0.041 & [\ensuremath{-}0.064, \ensuremath{-}0.021] & 0.000 / \ensuremath{-}0.083 & \ensuremath{-}0.009 & \ensuremath{-}0.003 & \ensuremath{+}0.002 & \ensuremath{+}0.026 \\
 &  & Macro-F1 & \ensuremath{-}0.002 & [\ensuremath{-}0.015, \ensuremath{+}0.011] & \ensuremath{-}0.001 / \ensuremath{-}0.002 & \ensuremath{+}0.001 & \ensuremath{+}0.002 & \ensuremath{+}0.001 & \ensuremath{+}0.005 \\
 &  & QWK+floor & 0.000 & [0.000, 0.000] & 0.000 / 0.000 & 0.000 & 0.000 & 0.000 & 0.000 \\
 &  & MMAE & \ensuremath{-}0.003 & [\ensuremath{-}0.013, \ensuremath{+}0.005] & 0.000 / \ensuremath{-}0.006 & \ensuremath{-}0.002 & 0.000 & 0.000 & \ensuremath{+}0.003 \\
 &  & MAE & \ensuremath{-}0.222 & [\ensuremath{-}0.277, \ensuremath{-}0.174] & \ensuremath{-}0.180 / \ensuremath{-}0.264 & \ensuremath{-}0.054 & \ensuremath{-}0.029 & 0.000 & \ensuremath{+}0.142 \\
 &  & RPS & \ensuremath{-}0.225 & [\ensuremath{-}0.277, \ensuremath{-}0.176] & \ensuremath{-}0.186 / \ensuremath{-}0.264 & \ensuremath{-}0.062 & \ensuremath{-}0.028 & \ensuremath{+}0.001 & \ensuremath{+}0.148 \\
\addlinespace
 & RF & AMAE & \ensuremath{-}0.050 & [\ensuremath{-}0.077, \ensuremath{-}0.025] & \ensuremath{-}0.037 / \ensuremath{-}0.063 & \ensuremath{-}0.013 & \ensuremath{-}0.005 & \ensuremath{+}0.001 & \ensuremath{+}0.015 \\
 &  & Macro-F1 & \ensuremath{-}0.050 & [\ensuremath{-}0.079, \ensuremath{-}0.025] & \ensuremath{-}0.037 / \ensuremath{-}0.063 & \ensuremath{-}0.013 & \ensuremath{-}0.005 & \ensuremath{+}0.001 & \ensuremath{+}0.015 \\
 &  & QWK+floor & 0.000 & [0.000, 0.000] & 0.000 / 0.000 & 0.000 & 0.000 & 0.000 & 0.000 \\
 &  & MMAE & \ensuremath{+}0.001 & [\ensuremath{-}0.020, \ensuremath{+}0.022] & 0.000 / \ensuremath{+}0.002 & \ensuremath{+}0.007 & \ensuremath{-}0.001 & \ensuremath{-}0.002 & \ensuremath{-}0.014 \\
 &  & MAE & \ensuremath{-}0.347 & [\ensuremath{-}0.401, \ensuremath{-}0.292] & \ensuremath{-}0.383 / \ensuremath{-}0.311 & \ensuremath{-}0.154 & \ensuremath{-}0.029 & \ensuremath{+}0.017 & \ensuremath{+}0.212 \\
 &  & RPS & \ensuremath{-}0.198 & [\ensuremath{-}0.244, \ensuremath{-}0.154] & \ensuremath{-}0.228 / \ensuremath{-}0.168 & \ensuremath{-}0.037 & \ensuremath{-}0.026 & 0.000 & \ensuremath{+}0.079 \\
\addlinespace
NPHA (control) & PO & AMAE & \ensuremath{-}0.298 & [\ensuremath{-}0.341, \ensuremath{-}0.256] & \ensuremath{-}0.354 / \ensuremath{-}0.241 & \ensuremath{-}0.012 & \ensuremath{-}0.134 & \ensuremath{-}0.071 & \ensuremath{+}0.066 \\
 &  & Macro-F1 & \ensuremath{+}0.069 & [\ensuremath{+}0.050, \ensuremath{+}0.091] & 0.000 / \ensuremath{+}0.138 & \ensuremath{+}0.014 & \ensuremath{+}0.016 & \ensuremath{+}0.003 & \ensuremath{-}0.004 \\
 &  & QWK+floor & \ensuremath{+}0.065 & [\ensuremath{+}0.047, \ensuremath{+}0.084] & 0.000 / \ensuremath{+}0.130 & \ensuremath{+}0.009 & \ensuremath{+}0.021 & \ensuremath{+}0.007 & \ensuremath{-}0.004 \\
 &  & MMAE & \ensuremath{+}0.067 & [\ensuremath{+}0.049, \ensuremath{+}0.087] & 0.000 / \ensuremath{+}0.134 & \ensuremath{+}0.013 & \ensuremath{+}0.022 & \ensuremath{+}0.007 & \ensuremath{-}0.006 \\
 &  & MAE & \ensuremath{-}0.304 & [\ensuremath{-}0.346, \ensuremath{-}0.263] & \ensuremath{-}0.354 / \ensuremath{-}0.253 & \ensuremath{-}0.020 & \ensuremath{-}0.134 & \ensuremath{-}0.070 & \ensuremath{+}0.064 \\
 &  & RPS & \ensuremath{-}0.304 & [\ensuremath{-}0.346, \ensuremath{-}0.262] & \ensuremath{-}0.354 / \ensuremath{-}0.253 & \ensuremath{-}0.020 & \ensuremath{-}0.134 & \ensuremath{-}0.070 & \ensuremath{+}0.064 \\
\addlinespace
 & RF & AMAE & \ensuremath{-}0.335 & [\ensuremath{-}0.382, \ensuremath{-}0.291] & \ensuremath{-}0.313 / \ensuremath{-}0.358 & \ensuremath{-}0.036 & \ensuremath{-}0.146 & \ensuremath{-}0.079 & \ensuremath{+}0.055 \\
 &  & Macro-F1 & \ensuremath{+}0.002 & [\ensuremath{-}0.021, \ensuremath{+}0.027] & \ensuremath{+}0.005 / \ensuremath{-}0.001 & \ensuremath{+}0.021 & \ensuremath{-}0.011 & \ensuremath{-}0.009 & \ensuremath{+}0.007 \\
 &  & QWK+floor & 0.000 & [0.000, 0.000] & 0.000 / 0.000 & 0.000 & 0.000 & 0.000 & 0.000 \\
 &  & MMAE & \ensuremath{-}0.036 & [\ensuremath{-}0.061, \ensuremath{-}0.011] & \ensuremath{+}0.048 / \ensuremath{-}0.121 & \ensuremath{-}0.001 & \ensuremath{-}0.013 & \ensuremath{-}0.007 & \ensuremath{+}0.018 \\
 &  & MAE & \ensuremath{-}0.335 & [\ensuremath{-}0.380, \ensuremath{-}0.292] & \ensuremath{-}0.313 / \ensuremath{-}0.358 & \ensuremath{-}0.036 & \ensuremath{-}0.146 & \ensuremath{-}0.079 & \ensuremath{+}0.055 \\
 &  & RPS & \ensuremath{-}0.335 & [\ensuremath{-}0.380, \ensuremath{-}0.291] & \ensuremath{-}0.313 / \ensuremath{-}0.358 & \ensuremath{-}0.036 & \ensuremath{-}0.146 & \ensuremath{-}0.079 & \ensuremath{+}0.055 \\
\addlinespace
SkillCraft, 8 leagues (superseded) & PO & AMAE & 0.000 & [0.000, 0.000] & 0.000 / 0.000 & 0.000 & 0.000 & 0.000 & 0.000 \\
 &  & Macro-F1 & \ensuremath{-}0.050 & [\ensuremath{-}0.073, \ensuremath{-}0.028] & \ensuremath{-}0.099 / 0.000 & \ensuremath{-}0.003 & \ensuremath{-}0.003 & \ensuremath{+}0.004 & \ensuremath{-}0.002 \\
 &  & QWK+floor & 0.000 & [0.000, 0.000] & 0.000 / 0.000 & 0.000 & 0.000 & 0.000 & 0.000 \\
 &  & MMAE & 0.000 & [0.000, 0.000] & 0.000 / 0.000 & 0.000 & 0.000 & 0.000 & 0.000 \\
 &  & MAE & \ensuremath{-}0.319 & [\ensuremath{-}0.389, \ensuremath{-}0.253] & \ensuremath{-}0.316 / \ensuremath{-}0.322 & \ensuremath{-}0.012 & \ensuremath{-}0.017 & \ensuremath{+}0.022 & \ensuremath{+}0.045 \\
 &  & RPS & \ensuremath{-}0.318 & [\ensuremath{-}0.389, \ensuremath{-}0.251] & \ensuremath{-}0.313 / \ensuremath{-}0.322 & \ensuremath{-}0.012 & \ensuremath{-}0.017 & \ensuremath{+}0.022 & \ensuremath{+}0.044 \\
\addlinespace
 & RF & AMAE & \ensuremath{+}0.043 & [\ensuremath{+}0.021, \ensuremath{+}0.067] & 0.000 / \ensuremath{+}0.087 & \ensuremath{-}0.013 & \ensuremath{+}0.010 & 0.000 & \ensuremath{+}0.015 \\
 &  & Macro-F1 & 0.000 & [0.000, 0.000] & 0.000 / 0.000 & 0.000 & 0.000 & 0.000 & 0.000 \\
 &  & QWK+floor & 0.000 & [0.000, 0.000] & 0.000 / 0.000 & 0.000 & 0.000 & 0.000 & 0.000 \\
 &  & MMAE & 0.000 & [0.000, 0.000] & 0.000 / 0.000 & 0.000 & 0.000 & 0.000 & 0.000 \\
 &  & MAE & \ensuremath{-}0.101 & [\ensuremath{-}0.142, \ensuremath{-}0.065] & \ensuremath{-}0.090 / \ensuremath{-}0.111 & \ensuremath{-}0.012 & \ensuremath{-}0.006 & \ensuremath{+}0.010 & \ensuremath{+}0.030 \\
 &  & RPS & \ensuremath{-}0.101 & [\ensuremath{-}0.140, \ensuremath{-}0.066] & \ensuremath{-}0.090 / \ensuremath{-}0.111 & \ensuremath{-}0.012 & \ensuremath{-}0.006 & \ensuremath{+}0.010 & \ensuremath{+}0.030 \\
\addlinespace

\end{longtable}
}

{\footnotesize
\begin{longtable}{@{}llrrrrrr@{}}
\caption{Fixed-hyperparameter controls: balanced minus unweighted training at the same family, hyperparameter, splits and seeds.}\label{tab:s-controls}\\
\toprule
Task & Configuration & $\Delta R_0$ & $\Delta R_{K-1}$ & $\Delta$QWK & $\Delta$MAE & $\Delta$AMAE & $\Delta$MMAE\\
\midrule\endfirsthead
\toprule
Task & Configuration & $\Delta R_0$ & $\Delta R_{K-1}$ & $\Delta$QWK & $\Delta$MAE & $\Delta$AMAE & $\Delta$MMAE\\
\midrule\endhead
\bottomrule\endfoot
Wine (red) & PO C=0.1 & \ensuremath{+}0.450 & \ensuremath{+}0.611 & \ensuremath{-}0.026 & \ensuremath{+}0.088 & \ensuremath{-}0.061 & \ensuremath{-}0.228 \\
 & PO C=1 & \ensuremath{+}0.400 & \ensuremath{+}0.611 & \ensuremath{-}0.040 & \ensuremath{+}0.089 & \ensuremath{-}0.055 & \ensuremath{-}0.211 \\
 & RF depth 4 & \ensuremath{+}0.050 & \ensuremath{+}0.306 & \ensuremath{+}0.012 & \ensuremath{+}0.060 & \ensuremath{-}0.064 & \ensuremath{-}0.170 \\
 & RF unlimited & 0.000 & \ensuremath{+}0.056 & \ensuremath{+}0.042 & \ensuremath{+}0.015 & \ensuremath{-}0.051 & \ensuremath{-}0.130 \\
Wine (white) & PO C=0.1 & \ensuremath{+}0.275 & \ensuremath{+}0.700 & \ensuremath{-}0.029 & \ensuremath{+}0.117 & \ensuremath{+}0.001 & \ensuremath{+}0.062 \\
 & PO C=1 & \ensuremath{+}0.300 & \ensuremath{+}0.700 & \ensuremath{-}0.038 & \ensuremath{+}0.117 & \ensuremath{+}0.009 & \ensuremath{+}0.071 \\
 & RF depth 4 & \ensuremath{+}0.125 & \ensuremath{+}0.100 & \ensuremath{+}0.031 & \ensuremath{+}0.073 & \ensuremath{-}0.036 & \ensuremath{-}0.096 \\
 & RF unlimited & \ensuremath{+}0.075 & 0.000 & \ensuremath{+}0.066 & \ensuremath{+}0.013 & \ensuremath{-}0.032 & \ensuremath{-}0.046 \\
SkillCraft & PO C=0.1 & \ensuremath{+}0.320 & \ensuremath{+}0.671 & \ensuremath{+}0.019 & \ensuremath{+}0.016 & \ensuremath{-}0.022 & \ensuremath{-}0.057 \\
 & PO C=1 & \ensuremath{+}0.320 & \ensuremath{+}0.671 & \ensuremath{+}0.020 & \ensuremath{+}0.014 & \ensuremath{-}0.023 & \ensuremath{-}0.056 \\
 & RF depth 4 & \ensuremath{+}0.353 & \ensuremath{+}0.586 & \ensuremath{+}0.012 & \ensuremath{+}0.027 & \ensuremath{-}0.018 & \ensuremath{-}0.038 \\
 & RF unlimited & \ensuremath{+}0.150 & \ensuremath{+}0.143 & \ensuremath{+}0.019 & \ensuremath{+}0.007 & \ensuremath{-}0.011 & \ensuremath{-}0.039 \\
Restaurant & PO C=0.1 & \ensuremath{+}0.352 & \ensuremath{-}0.207 & \ensuremath{-}0.029 & \ensuremath{+}0.100 & \ensuremath{+}0.044 & \ensuremath{-}0.255 \\
 & PO C=1 & \ensuremath{+}0.348 & \ensuremath{-}0.183 & \ensuremath{-}0.025 & \ensuremath{+}0.091 & \ensuremath{+}0.039 & \ensuremath{-}0.237 \\
 & RF depth 4 & \ensuremath{+}0.317 & \ensuremath{-}0.352 & \ensuremath{-}0.043 & \ensuremath{+}0.083 & \ensuremath{+}0.018 & \ensuremath{-}0.298 \\
 & RF unlimited & \ensuremath{+}0.213 & \ensuremath{-}0.114 & \ensuremath{+}0.047 & \ensuremath{+}0.022 & \ensuremath{-}0.012 & \ensuremath{-}0.192 \\
Sushi & PO C=0.1 & \ensuremath{+}0.458 & \ensuremath{-}0.382 & \ensuremath{+}0.072 & \ensuremath{+}0.092 & \ensuremath{-}0.044 & \ensuremath{-}0.419 \\
 & PO C=1 & \ensuremath{+}0.450 & \ensuremath{-}0.355 & \ensuremath{+}0.060 & \ensuremath{+}0.094 & \ensuremath{-}0.037 & \ensuremath{-}0.405 \\
 & RF depth 4 & \ensuremath{+}0.318 & \ensuremath{-}0.613 & \ensuremath{+}0.139 & \ensuremath{+}0.039 & \ensuremath{-}0.134 & \ensuremath{-}0.527 \\
 & RF unlimited & \ensuremath{+}0.338 & \ensuremath{-}0.424 & \ensuremath{+}0.066 & \ensuremath{+}0.082 & \ensuremath{-}0.061 & \ensuremath{-}0.386 \\
MovieLens-100k & PO C=0.1 & \ensuremath{+}0.519 & \ensuremath{+}0.288 & \ensuremath{+}0.060 & \ensuremath{+}0.149 & \ensuremath{+}0.031 & \ensuremath{-}0.281 \\
 & PO C=1 & \ensuremath{+}0.519 & \ensuremath{+}0.288 & \ensuremath{+}0.060 & \ensuremath{+}0.149 & \ensuremath{+}0.031 & \ensuremath{-}0.281 \\
 & RF depth 4 & \ensuremath{+}0.426 & \ensuremath{+}0.530 & \ensuremath{+}0.172 & \ensuremath{+}0.152 & \ensuremath{+}0.021 & \ensuremath{-}0.343 \\
 & RF unlimited & \ensuremath{+}0.205 & \ensuremath{+}0.083 & \ensuremath{+}0.017 & \ensuremath{+}0.068 & \ensuremath{+}0.006 & \ensuremath{-}0.198 \\
CTG & PO C=0.1 & \ensuremath{-}0.147 & \ensuremath{+}0.219 & \ensuremath{-}0.031 & \ensuremath{+}0.030 & \ensuremath{-}0.068 & \ensuremath{-}0.139 \\
 & PO C=1 & \ensuremath{-}0.129 & \ensuremath{+}0.224 & \ensuremath{-}0.016 & \ensuremath{+}0.025 & \ensuremath{-}0.068 & \ensuremath{-}0.152 \\
 & RF depth 4 & \ensuremath{-}0.107 & \ensuremath{+}0.168 & \ensuremath{-}0.030 & \ensuremath{+}0.020 & \ensuremath{-}0.076 & \ensuremath{-}0.152 \\
 & RF unlimited & \ensuremath{-}0.049 & \ensuremath{+}0.088 & \ensuremath{-}0.034 & \ensuremath{+}0.011 & \ensuremath{-}0.030 & \ensuremath{-}0.062 \\
HCV & PO C=0.1 & \ensuremath{-}0.337 & \ensuremath{+}0.050 & \ensuremath{-}0.053 & \ensuremath{+}0.079 & \ensuremath{-}0.152 & \ensuremath{-}0.322 \\
 & PO C=1 & \ensuremath{-}0.322 & \ensuremath{+}0.033 & \ensuremath{-}0.111 & \ensuremath{+}0.082 & \ensuremath{-}0.112 & \ensuremath{-}0.211 \\
 & RF depth 4 & \ensuremath{-}0.058 & \ensuremath{+}0.067 & \ensuremath{+}0.046 & \ensuremath{+}0.006 & \ensuremath{-}0.128 & \ensuremath{-}0.244 \\
 & RF unlimited & \ensuremath{-}0.014 & \ensuremath{+}0.100 & \ensuremath{+}0.063 & \ensuremath{-}0.006 & \ensuremath{-}0.082 & \ensuremath{-}0.103 \\
Cirrhosis & PO C=0.1 & \ensuremath{+}0.405 & \ensuremath{-}0.028 & \ensuremath{+}0.008 & \ensuremath{+}0.080 & \ensuremath{-}0.003 & \ensuremath{-}0.256 \\
 & PO C=1 & \ensuremath{+}0.429 & \ensuremath{-}0.031 & \ensuremath{-}0.012 & \ensuremath{+}0.079 & \ensuremath{+}0.005 & \ensuremath{-}0.222 \\
 & RF depth 4 & \ensuremath{+}0.452 & \ensuremath{+}0.056 & \ensuremath{+}0.113 & \ensuremath{+}0.038 & \ensuremath{-}0.067 & \ensuremath{-}0.381 \\
 & RF unlimited & \ensuremath{+}0.238 & \ensuremath{-}0.010 & \ensuremath{+}0.054 & \ensuremath{+}0.022 & \ensuremath{-}0.036 & \ensuremath{-}0.214 \\
SUPPORT2 & PO C=0.1 & \ensuremath{-}0.213 & \ensuremath{-}0.346 & \ensuremath{+}0.006 & \ensuremath{+}0.009 & \ensuremath{-}0.072 & \ensuremath{-}0.128 \\
 & PO C=1 & \ensuremath{-}0.216 & \ensuremath{-}0.344 & \ensuremath{+}0.004 & \ensuremath{+}0.010 & \ensuremath{-}0.070 & \ensuremath{-}0.129 \\
 & RF depth 4 & \ensuremath{-}0.150 & \ensuremath{-}0.357 & \ensuremath{-}0.078 & \ensuremath{+}0.022 & \ensuremath{-}0.070 & \ensuremath{-}0.061 \\
 & RF unlimited & \ensuremath{-}0.051 & \ensuremath{-}0.102 & \ensuremath{-}0.018 & \ensuremath{+}0.002 & \ensuremath{-}0.028 & \ensuremath{-}0.138 \\
TAE (control) & PO C=0.1 & \ensuremath{+}0.204 & \ensuremath{-}0.019 & \ensuremath{+}0.054 & \ensuremath{-}0.003 & \ensuremath{-}0.004 & \ensuremath{+}0.019 \\
 & PO C=1 & \ensuremath{+}0.061 & 0.000 & \ensuremath{+}0.026 & \ensuremath{-}0.013 & \ensuremath{-}0.014 & 0.000 \\
 & RF depth 4 & 0.000 & \ensuremath{+}0.058 & \ensuremath{+}0.063 & \ensuremath{-}0.020 & \ensuremath{-}0.019 & \ensuremath{-}0.048 \\
 & RF unlimited & \ensuremath{-}0.041 & \ensuremath{+}0.010 & \ensuremath{-}0.022 & \ensuremath{+}0.007 & \ensuremath{+}0.007 & \ensuremath{-}0.005 \\
Diabetes130 (control) & PO C=0.1 & \ensuremath{-}0.241 & \ensuremath{+}0.396 & \ensuremath{+}0.049 & \ensuremath{+}0.067 & \ensuremath{-}0.037 & \ensuremath{-}0.327 \\
 & PO C=1 & \ensuremath{-}0.232 & \ensuremath{+}0.377 & \ensuremath{+}0.039 & \ensuremath{+}0.068 & \ensuremath{-}0.031 & \ensuremath{-}0.304 \\
 & RF depth 4 & \ensuremath{-}0.316 & \ensuremath{+}0.463 & \ensuremath{+}0.163 & \ensuremath{+}0.058 & \ensuremath{-}0.091 & \ensuremath{-}0.493 \\
 & RF unlimited & \ensuremath{-}0.194 & \ensuremath{+}0.192 & \ensuremath{+}0.050 & \ensuremath{+}0.026 & \ensuremath{-}0.039 & \ensuremath{-}0.204 \\
Student (control) & PO C=0.1 & \ensuremath{+}0.225 & \ensuremath{+}0.226 & \ensuremath{+}0.062 & \ensuremath{+}0.028 & \ensuremath{-}0.001 & \ensuremath{-}0.148 \\
 & PO C=1 & \ensuremath{+}0.230 & \ensuremath{+}0.189 & \ensuremath{+}0.046 & \ensuremath{+}0.023 & \ensuremath{-}0.003 & \ensuremath{-}0.125 \\
 & RF depth 4 & \ensuremath{+}0.545 & \ensuremath{+}0.268 & \ensuremath{+}0.179 & \ensuremath{+}0.037 & \ensuremath{-}0.018 & \ensuremath{-}0.258 \\
 & RF unlimited & \ensuremath{+}0.305 & \ensuremath{+}0.140 & \ensuremath{+}0.067 & \ensuremath{+}0.020 & \ensuremath{-}0.009 & \ensuremath{-}0.128 \\
NPHA (control) & PO C=0.1 & \ensuremath{+}0.485 & \ensuremath{+}0.318 & \ensuremath{+}0.047 & \ensuremath{+}0.161 & \ensuremath{+}0.076 & \ensuremath{-}0.071 \\
 & PO C=1 & \ensuremath{+}0.431 & \ensuremath{+}0.244 & \ensuremath{+}0.023 & \ensuremath{+}0.140 & \ensuremath{+}0.068 & \ensuremath{-}0.078 \\
 & RF depth 4 & \ensuremath{+}0.374 & \ensuremath{+}0.360 & \ensuremath{+}0.054 & \ensuremath{+}0.155 & \ensuremath{+}0.081 & \ensuremath{-}0.063 \\
 & RF unlimited & \ensuremath{+}0.286 & \ensuremath{+}0.225 & \ensuremath{+}0.048 & \ensuremath{+}0.100 & \ensuremath{+}0.049 & \ensuremath{-}0.071 \\
SkillCraft, 8 leagues (superseded) & PO C=0.1 & \ensuremath{+}0.347 & \ensuremath{+}0.291 & \ensuremath{+}0.013 & \ensuremath{+}0.017 & \ensuremath{-}0.022 & \ensuremath{-}0.046 \\
 & PO C=1 & \ensuremath{+}0.341 & \ensuremath{+}0.291 & \ensuremath{+}0.012 & \ensuremath{+}0.017 & \ensuremath{-}0.022 & \ensuremath{-}0.044 \\
 & RF depth 4 & \ensuremath{+}0.395 & \ensuremath{+}0.345 & \ensuremath{+}0.014 & \ensuremath{+}0.024 & \ensuremath{-}0.027 & \ensuremath{-}0.033 \\
 & RF unlimited & \ensuremath{+}0.129 & \ensuremath{+}0.073 & \ensuremath{+}0.012 & \ensuremath{+}0.006 & \ensuremath{-}0.010 & \ensuremath{-}0.030 \\

\end{longtable}
}

{\footnotesize
\begin{longtable}{@{}lllllrlrr@{}}
\caption{Controlled simulation: differences relative to QWK selection under the latent Gaussian model (ten replications per cell; Monte Carlo 95\% $t$ intervals, 9 df). Same: fraction of replications in which the two criteria selected the same candidate.}\label{tab:s-sim}\\
\toprule
Prior & Signal & Fam. & Criterion & Same & $\Delta R_{\mathrm{end}}$ & MC 95\% interval & $\Delta$QWK & $\Delta$AMAE\\
\midrule\endfirsthead
\toprule
Prior & Signal & Fam. & Criterion & Same & $\Delta R_{\mathrm{end}}$ & MC 95\% interval & $\Delta$QWK & $\Delta$AMAE\\
\midrule\endhead
\bottomrule\endfoot
balanced & 0.5 & PO & AMAE & 0.70 & \ensuremath{-}0.002 & [\ensuremath{-}0.005, \ensuremath{+}0.002] & 0.000 & \ensuremath{-}0.001 \\
balanced & 0.5 & PO & Macro-F1 & 0.50 & \ensuremath{+}0.006 & [\ensuremath{-}0.008, \ensuremath{+}0.021] & 0.000 & \ensuremath{+}0.001 \\
balanced & 0.5 & RF & AMAE & 0.60 & \ensuremath{-}0.058 & [\ensuremath{-}0.113, \ensuremath{-}0.004] & \ensuremath{-}0.018 & \ensuremath{-}0.002 \\
balanced & 0.5 & RF & Macro-F1 & 0.60 & \ensuremath{-}0.054 & [\ensuremath{-}0.105, \ensuremath{-}0.002] & \ensuremath{-}0.020 & 0.000 \\
balanced & 2.0 & PO & AMAE & 0.90 & \ensuremath{+}0.001 & [\ensuremath{-}0.001, \ensuremath{+}0.002] & 0.000 & 0.000 \\
balanced & 2.0 & PO & Macro-F1 & 0.80 & \ensuremath{+}0.001 & [\ensuremath{-}0.001, \ensuremath{+}0.002] & 0.000 & 0.000 \\
balanced & 2.0 & RF & AMAE & 0.70 & \ensuremath{+}0.001 & [\ensuremath{-}0.007, \ensuremath{+}0.009] & \ensuremath{-}0.001 & 0.000 \\
balanced & 2.0 & RF & Macro-F1 & 0.50 & \ensuremath{+}0.007 & [\ensuremath{-}0.006, \ensuremath{+}0.021] & \ensuremath{-}0.002 & \ensuremath{+}0.001 \\
symmetric rare & 0.5 & PO & AMAE & 1.00 & 0.000 & [0.000, 0.000] & 0.000 & 0.000 \\
symmetric rare & 0.5 & PO & Macro-F1 & 0.80 & 0.000 & [0.000, 0.000] & 0.000 & 0.000 \\
symmetric rare & 0.5 & RF & AMAE & 0.60 & \ensuremath{-}0.077 & [\ensuremath{-}0.201, \ensuremath{+}0.048] & \ensuremath{-}0.002 & 0.000 \\
symmetric rare & 0.5 & RF & Macro-F1 & 0.40 & \ensuremath{-}0.171 & [\ensuremath{-}0.280, \ensuremath{-}0.062] & \ensuremath{-}0.001 & 0.000 \\
symmetric rare & 2.0 & PO & AMAE & 0.90 & \ensuremath{-}0.001 & [\ensuremath{-}0.003, \ensuremath{+}0.001] & 0.000 & 0.000 \\
symmetric rare & 2.0 & PO & Macro-F1 & 0.40 & \ensuremath{-}0.146 & [\ensuremath{-}0.258, \ensuremath{-}0.034] & \ensuremath{-}0.004 & \ensuremath{+}0.015 \\
symmetric rare & 2.0 & RF & AMAE & 0.60 & \ensuremath{+}0.072 & [\ensuremath{-}0.003, \ensuremath{+}0.146] & \ensuremath{-}0.011 & \ensuremath{-}0.002 \\
symmetric rare & 2.0 & RF & Macro-F1 & 1.00 & 0.000 & [0.000, 0.000] & 0.000 & 0.000 \\
asymmetric & 0.5 & PO & AMAE & 0.60 & 0.000 & [\ensuremath{-}0.001, 0.000] & 0.000 & 0.000 \\
asymmetric & 0.5 & PO & Macro-F1 & 0.40 & 0.000 & [\ensuremath{-}0.001, \ensuremath{+}0.001] & 0.000 & 0.000 \\
asymmetric & 0.5 & RF & AMAE & 0.80 & \ensuremath{+}0.017 & [\ensuremath{-}0.009, \ensuremath{+}0.042] & \ensuremath{+}0.004 & \ensuremath{-}0.004 \\
asymmetric & 0.5 & RF & Macro-F1 & 0.60 & \ensuremath{-}0.058 & [\ensuremath{-}0.113, \ensuremath{-}0.004] & \ensuremath{-}0.012 & \ensuremath{+}0.004 \\
asymmetric & 2.0 & PO & AMAE & 0.40 & \ensuremath{+}0.020 & [\ensuremath{+}0.004, \ensuremath{+}0.036] & \ensuremath{-}0.003 & \ensuremath{-}0.006 \\
asymmetric & 2.0 & PO & Macro-F1 & 0.50 & \ensuremath{-}0.007 & [\ensuremath{-}0.027, \ensuremath{+}0.013] & \ensuremath{-}0.001 & \ensuremath{+}0.002 \\
asymmetric & 2.0 & RF & AMAE & 0.60 & 0.000 & [\ensuremath{-}0.023, \ensuremath{+}0.023] & \ensuremath{-}0.003 & \ensuremath{-}0.001 \\
asymmetric & 2.0 & RF & Macro-F1 & 0.70 & \ensuremath{-}0.016 & [\ensuremath{-}0.037, \ensuremath{+}0.005] & 0.000 & \ensuremath{+}0.001 \\

\end{longtable}
}

{\footnotesize
\begin{longtable}{@{}lrr*{7}{r}@{}}
\caption{Friedman tests and mean ranks of the seven criteria for each evaluation metric, with the \FriedBlocks{} primary sources as blocks: ranks are computed within each of the \NCells{} task--family cells and averaged within each source (the two model families, and for wine the two tasks), because these cells are not independent. Rank 1 = best; Nemenyi critical difference at $\alpha=0.05$ with \FriedBlocks{} blocks: \NemenyiCD{}.}\label{tab:s-friedman}\\
\toprule
Evaluation & $\chi^2_6$ & $p$ & QWK & MAE & AMAE & MMAE & Macro-F1 & RPS & QWK-floor\\
\midrule\endfirsthead
\bottomrule\endfoot
QWK & 8.5 & 0.204 & 2.44 & 4.39 & 4.44 & 4.89 & 4.00 & 4.50 & 3.33 \\
Macro-F1 & 19.9 & 0.003 & 2.89 & 5.61 & 3.39 & 3.39 & 2.89 & 6.00 & 3.83 \\
MAE & 33.3 & \ensuremath{<}0.001 & 3.67 & 1.89 & 5.11 & 5.67 & 4.89 & 1.56 & 5.22 \\
AMAE & 18.0 & 0.006 & 3.72 & 5.56 & 2.17 & 3.28 & 3.94 & 5.61 & 3.72 \\
MMAE & 40.8 & \ensuremath{<}0.001 & 4.00 & 6.33 & 2.72 & 1.67 & 3.44 & 6.67 & 3.17 \\
RPS & 37.2 & \ensuremath{<}0.001 & 3.94 & 1.94 & 4.94 & 6.11 & 5.06 & 1.28 & 4.72 \\
Endpoint recall & 13.3 & 0.039 & 3.89 & 5.39 & 4.56 & 2.89 & 3.94 & 4.89 & 2.44 \\

\end{longtable}
}

{\footnotesize
\begin{longtable}{@{}l*{7}{r}@{}}
\caption{Stage 4, expanded pool: fraction of decisions selecting class-balanced weights, by family and criterion (primary tasks; PO: proportional odds, RF: random forest, HGB: histogram gradient boosting).}\label{tab:s-stage4-balance}\\
\toprule
Family & QWK & MAE & AMAE & MMAE & Macro-F1 & RPS & QWK-floor\\
\midrule\endfirsthead
\bottomrule\endfoot
PO & 0.45 & 0.00 & 0.73 & 0.93 & 0.57 & 0.00 & 0.75\\
RF & 0.75 & 0.07 & 0.88 & 1.00 & 0.88 & 0.00 & 0.80\\
HGB & 0.62 & 0.05 & 0.82 & 0.85 & 0.73 & 0.02 & 0.68\\

\end{longtable}
}

{\footnotesize
\begin{longtable}{@{}l*{7}{r}@{}}
\caption{Stage 4, expanded pool: normalized cross-criterion regret (rows: selection criterion; columns: evaluation metric), averaged over the primary task--family cells and two repeats; 0 is the best and 1 the worst of the eight candidates in the family, identified retrospectively.}\label{tab:s-stage4-regret}\\
\toprule
Criterion & QWK & Macro-F1 & MAE & AMAE & MMAE & RPS & $R_{\mathrm{end}}$\\
\midrule\endfirsthead
\bottomrule\endfoot
QWK & 0.15 & 0.18 & 0.48 & 0.37 & 0.32 & 0.44 & 0.43\\
MAE & 0.46 & 0.55 & 0.08 & 0.58 & 0.74 & 0.09 & 0.65\\
AMAE & 0.34 & 0.31 & 0.63 & 0.13 & 0.24 & 0.56 & 0.55\\
MMAE & 0.33 & 0.31 & 0.76 & 0.28 & 0.08 & 0.68 & 0.40\\
Macro-F1 & 0.27 & 0.19 & 0.54 & 0.32 & 0.28 & 0.54 & 0.54\\
RPS & 0.40 & 0.55 & 0.06 & 0.58 & 0.77 & 0.02 & 0.61\\
QWK-floor & 0.23 & 0.26 & 0.62 & 0.35 & 0.24 & 0.55 & 0.34\\

\end{longtable}
}

{\footnotesize
\begin{longtable}{@{}rrrllll@{}}
\caption{Stage 4 simulation: change in the recall of the lowest category. $\Delta_w$: balanced minus unweighted training (mean of the two hyperparameter settings) with Monte Carlo 95\% $t$ interval (9 df); $\Delta_a$: AMAE-selected minus QWK-selected model, with the number of replications (of ten) in which it was nonzero.}\label{tab:s-stage4-sim}\\
\toprule
& & & \multicolumn{2}{c}{Proportional odds} & \multicolumn{2}{c}{Random forest}\\
$\lambda$ & $\pi_0$ & $K\pi_0$ & $\Delta_w$ [95\% interval] & $\Delta_a$ & $\Delta_w$ [95\% interval] & $\Delta_a$\\
\midrule\endfirsthead
\toprule
$\lambda$ & $\pi_0$ & $K\pi_0$ & $\Delta_w$ [95\% interval] & $\Delta_a$ & $\Delta_w$ [95\% interval] & $\Delta_a$\\
\midrule\endhead
\bottomrule\endfoot
0.5 & 0.02 & 0.10 & \ensuremath{+}0.543 [\ensuremath{+}0.477, \ensuremath{+}0.609] & \ensuremath{-}0.100 (2) & \ensuremath{+}0.224 [\ensuremath{+}0.203, \ensuremath{+}0.245] & \ensuremath{-}0.079 (6)\\
0.5 & 0.05 & 0.25 & \ensuremath{+}0.503 [\ensuremath{+}0.473, \ensuremath{+}0.533] & \ensuremath{-}0.246 (5) & \ensuremath{+}0.274 [\ensuremath{+}0.238, \ensuremath{+}0.309] & \ensuremath{-}0.246 (6)\\
0.5 & 0.10 & 0.50 & \ensuremath{+}0.478 [\ensuremath{+}0.446, \ensuremath{+}0.511] & \ensuremath{-}0.342 (8) & \ensuremath{+}0.296 [\ensuremath{+}0.263, \ensuremath{+}0.329] & \ensuremath{-}0.297 (7)\\
0.5 & 0.20 & 1.00 & \ensuremath{-}0.015 [\ensuremath{-}0.068, \ensuremath{+}0.039] & \ensuremath{-}0.004 (3) & \ensuremath{-}0.006 [\ensuremath{-}0.029, \ensuremath{+}0.017] & \ensuremath{-}0.064 (5)\\
0.5 & 0.30 & 1.50 & \ensuremath{-}0.306 [\ensuremath{-}0.331, \ensuremath{-}0.282] & 0.000 (1) & \ensuremath{-}0.268 [\ensuremath{-}0.295, \ensuremath{-}0.241] & \ensuremath{-}0.087 (4)\\
0.5 & 0.40 & 2.00 & \ensuremath{-}0.348 [\ensuremath{-}0.372, \ensuremath{-}0.324] & 0.000 (3) & \ensuremath{-}0.350 [\ensuremath{-}0.376, \ensuremath{-}0.324] & \ensuremath{-}0.044 (2)\\
1 & 0.02 & 0.10 & \ensuremath{+}0.629 [\ensuremath{+}0.572, \ensuremath{+}0.687] & \ensuremath{+}0.330 (5) & \ensuremath{+}0.324 [\ensuremath{+}0.256, \ensuremath{+}0.393] & \ensuremath{+}0.198 (5)\\
1 & 0.05 & 0.25 & \ensuremath{+}0.517 [\ensuremath{+}0.479, \ensuremath{+}0.555] & \ensuremath{+}0.052 (4) & \ensuremath{+}0.406 [\ensuremath{+}0.365, \ensuremath{+}0.447] & \ensuremath{+}0.056 (5)\\
1 & 0.10 & 0.50 & \ensuremath{+}0.343 [\ensuremath{+}0.303, \ensuremath{+}0.383] & \ensuremath{-}0.080 (5) & \ensuremath{+}0.365 [\ensuremath{+}0.323, \ensuremath{+}0.407] & \ensuremath{-}0.136 (4)\\
1 & 0.20 & 1.00 & \ensuremath{-}0.007 [\ensuremath{-}0.044, \ensuremath{+}0.030] & \ensuremath{-}0.011 (2) & \ensuremath{-}0.008 [\ensuremath{-}0.045, \ensuremath{+}0.029] & \ensuremath{-}0.024 (3)\\
1 & 0.30 & 1.50 & \ensuremath{-}0.250 [\ensuremath{-}0.269, \ensuremath{-}0.232] & 0.000 (1) & \ensuremath{-}0.192 [\ensuremath{-}0.209, \ensuremath{-}0.176] & \ensuremath{-}0.001 (2)\\
1 & 0.40 & 2.00 & \ensuremath{-}0.342 [\ensuremath{-}0.358, \ensuremath{-}0.326] & 0.000 (2) & \ensuremath{-}0.260 [\ensuremath{-}0.275, \ensuremath{-}0.246] & 0.000 (0)\\
2 & 0.02 & 0.10 & \ensuremath{+}0.505 [\ensuremath{+}0.459, \ensuremath{+}0.551] & \ensuremath{+}0.242 (6) & \ensuremath{+}0.461 [\ensuremath{+}0.388, \ensuremath{+}0.533] & \ensuremath{+}0.195 (5)\\
2 & 0.05 & 0.25 & \ensuremath{+}0.350 [\ensuremath{+}0.316, \ensuremath{+}0.383] & \ensuremath{+}0.057 (2) & \ensuremath{+}0.490 [\ensuremath{+}0.455, \ensuremath{+}0.525] & \ensuremath{+}0.108 (4)\\
2 & 0.10 & 0.50 & \ensuremath{+}0.187 [\ensuremath{+}0.155, \ensuremath{+}0.219] & \ensuremath{+}0.027 (5) & \ensuremath{+}0.306 [\ensuremath{+}0.241, \ensuremath{+}0.371] & \ensuremath{+}0.038 (2)\\
2 & 0.20 & 1.00 & \ensuremath{+}0.008 [\ensuremath{-}0.012, \ensuremath{+}0.028] & 0.000 (6) & \ensuremath{+}0.011 [\ensuremath{-}0.022, \ensuremath{+}0.045] & \ensuremath{-}0.019 (4)\\
2 & 0.30 & 1.50 & \ensuremath{-}0.123 [\ensuremath{-}0.143, \ensuremath{-}0.103] & \ensuremath{-}0.056 (6) & \ensuremath{-}0.127 [\ensuremath{-}0.151, \ensuremath{-}0.104] & \ensuremath{-}0.026 (4)\\
2 & 0.40 & 2.00 & \ensuremath{-}0.227 [\ensuremath{-}0.241, \ensuremath{-}0.212] & \ensuremath{-}0.022 (5) & \ensuremath{-}0.184 [\ensuremath{-}0.199, \ensuremath{-}0.169] & \ensuremath{-}0.067 (6)\\

\end{longtable}
}

{\footnotesize
\begin{longtable}{@{}lrrll@{}}
\caption{Stage 5: new primary tasks. $K\pi$: share of the lowest/highest category times $K$; best QWK: highest outer QWK of the 32 candidates (mean of two repeats).}\label{tab:s-stage5-tasks}\\
\toprule
Task & $n$ & $K$ & $K\pi$ (low / high) & Best QWK\\
\midrule\endfirsthead
\bottomrule\endfoot
Drug consumption (alcohol) & 1,885 & 7 & 0.13 / 1.88 & 0.15\\
Entrance examination & 666 & 4 & 0.94 / 0.61 & 0.62\\
NHANES self-rated health & 4,890 & 5 & 0.43 / 0.16 & 0.38\\
NHIS self-rated health & 5,000 & 5 & 1.12 / 0.19 & 0.55\\

\end{longtable}
}

{\footnotesize
\begin{longtable}{@{}lrrllrll@{}}
\caption{Endpoint-share rule with task-cluster and source-cluster bootstrap intervals. Agreement: share of endpoint--family points whose recall changed in the direction given by the rule (zero changes excluded for AMAE$-$QWK selection); $\rho$: Spearman correlation between the change and $\log(K\pi_c)$. The 95\% percentile intervals resample tasks (as in the main text) or the \NSourcesRule{} primary sources (the two wine tasks form one source), 2000 resamples each.}\label{tab:s-source-rule}\\
\toprule
& & \multicolumn{3}{c}{Agreement} & \multicolumn{3}{c}{Spearman $\rho$}\\
Comparison & Points & Estimate & Task clusters & Source clusters & Estimate & Task clusters & Source clusters\\
\midrule\endfirsthead
\bottomrule\endfoot
Stages 1--3: balanced $-$ unweighted & 40 & 0.93 & [0.80, 1.00] & [0.81, 1.00] & \ensuremath{-}0.66 & [\ensuremath{-}0.81, \ensuremath{-}0.37] & [\ensuremath{-}0.86, \ensuremath{-}0.32]\\
Stages 1--3: AMAE $-$ QWK selection & 40 & 0.77 & [0.63, 0.89] & [0.62, 0.90] & \ensuremath{-}0.67 & [\ensuremath{-}0.82, \ensuremath{-}0.32] & [\ensuremath{-}0.80, \ensuremath{-}0.23]\\
Stage 4 (H2): balanced $-$ unweighted & 60 & 0.90 & [0.80, 0.98] & [0.79, 0.98] & \ensuremath{-}0.60 & [\ensuremath{-}0.75, \ensuremath{-}0.26] & [\ensuremath{-}0.77, \ensuremath{-}0.25]\\
Stage 4 (H3): AMAE $-$ QWK selection & 60 & 0.81 & [0.67, 0.93] & [0.67, 0.94] & \ensuremath{-}0.62 & [\ensuremath{-}0.78, \ensuremath{-}0.30] & [\ensuremath{-}0.76, \ensuremath{-}0.19]\\

\end{longtable}
}

\end{document}
